\subsection{Duality of extension modules and torsion products}

Let A be a ring, and let P and J be A-modules. For every complex C of A-modules, we constructed in A, X, p.~99, prop. 12 a canonical isomorphism of complexes

\[
\mu : \operatorname{Homgr} _ {\mathrm{A}} (\mathrm{C} \otimes_ {\mathrm{A}} \mathrm{P}, \mathrm{J}) \longrightarrow \operatorname{Homgr} _ {\mathrm{A}} (\mathrm{C}, \operatorname{Hom} _ {\mathrm{A}} (\mathrm{P}, \mathrm{J})) .
\]

Let M be a \(\Lambda\) module, and (C, p) a projective resolution of M. Consider the sequence of homomorphisms

\[
\begin{array}{r l} \operatorname{Ext} _ {\mathrm{A}} (\mathrm{M}, \operatorname{Hom} _ {\mathrm{A}} (\mathrm{P}, \mathrm{J})) & \xrightarrow {\varphi^ {- i}} \mathrm{H} (\operatorname{Homgr} _ {\mathrm{A}} (\mathrm{C}, \operatorname{Hom} _ {\mathrm{A}} (\mathrm{P}, \mathrm{J}))) \xrightarrow {\mathrm{H} (\mu) ^ {- 1}} \mathrm{H} (\operatorname{Homgr} _ {\mathrm{A}} (\mathrm{C} \otimes_ {\mathrm{A}} \mathrm{P}, \mathrm{J})) \\ & \xrightarrow {u} \operatorname{Homgr} _ {\mathrm{A}} (\mathrm{H} (\mathrm{C} \otimes_ {\mathrm{A}} \mathrm{P}), \mathrm{J}) \xrightarrow {v} \operatorname{Homgr} _ {\mathrm{A}} (\operatorname{Tor} ^ {\mathrm{A}} (\mathrm{M}, \mathrm{P}), \mathrm{J}), \end{array}
\]

where \(\varphi\) is the canonical isomorphism \(\varphi(\mathrm{C},\mathrm{Hom}_{\mathrm{A}}(\mathrm{P},\mathrm{J}))\) (A, X, p.~100, th. 1) \(u\) the canonical homomorphism \(\lambda(\mathrm{C}\otimes_{\Lambda}\mathrm{P},\mathrm{J})\) . (A, X, p.~82), and \(v\) is deduced from the canonical isomorphism \(\psi(\mathrm{C},\mathrm{P}):\operatorname{Tor}^{\Lambda}(\mathrm{M},\mathrm{P})\longrightarrow\mathrm{H}(\mathrm{C}\otimes_{\mathrm{A}}\mathrm{P})\) .

Let \((\mathbf{C}', p')\) another projective resolution of M. By A, X, p.~49, cor. of prop. 3, there exists a homotopy of complexes \(\alpha: \mathbf{C}' \to \mathbf{C}\) such that \(p \circ \alpha = p'\) . It follows from A, X, p.~103, prop. 2, that one has \(\mathrm{H}(\alpha \otimes 1_{\mathrm{P}}) \circ \psi(\mathrm{C}', \mathrm{P}) = \psi(\mathrm{C}, \mathrm{P})\) and \(\varphi(\mathrm{C}', \mathrm{R}) \circ \mathrm{H}(\mathrm{Homgr}(\alpha, 1_{\mathrm{R}})) = \varphi(\mathrm{C}, \mathrm{R})\) for every A-module R. From this it follows that the graded homomorphism of degree 0

\[
\theta (\mathrm{M}, \mathrm{P}): \operatorname{Ext} _ {\mathrm{A}} (\mathrm{M}, \operatorname{Hom} _ {\mathrm{A}} (\mathrm{P}, \mathrm{J})) \longrightarrow \operatorname{Homgr} _ {\mathrm{A}} \left(\operatorname{Tor} ^ {\mathrm{A}} (\mathrm{M}, \mathrm{P}), \mathrm{J}\right)
\]

the composite of the sequence of homomorphisms above is independent of the choice of the projective resolution (C,p) of M. By construction it is End \(_{A}\) (J)-linear.

The definition of the homomorphism \(\theta(\mathrm{M},\mathrm{P})\) is made explicit as follows. Let \(p\) an integer, \(\mathfrak{v}\) be an element of \(\operatorname{Ext}_{\Lambda}^{p}(\mathrm{M},\operatorname{Hom}_{\mathrm{A}}(\mathrm{P},\mathrm{J}))\) , \(\tau\) be an element of \(\operatorname{Tor}_{p}^{\mathrm{A}}(\mathrm{M},\mathrm{P})\) . Using the isomorphism \(\varphi(\mathrm{C},\operatorname{Hom}_{\mathrm{A}}(\mathrm{P},\mathrm{J}))\) , \(\mathfrak{v}\) is represented by a linear map \(u:\mathrm{C}_{p}\to\operatorname{Hom}_{\mathrm{A}}(\mathrm{P},\mathrm{J})\) such that \(u\circ d_{\mathrm{C}}=0\) ; similarly, using \(\psi(\mathrm{C},\mathrm{P})\) , \(\tau\) is represented by an element \(\sum c_{\mu}\otimes p_{\mu}\) of \(\mathrm{C}_{p}\otimes\mathrm{P}\) such that \(\sum d_{\mathrm{C}}(c_{\mu})\otimes p_{\mu}=0\) . One then has \(\theta(\mathrm{M},\mathrm{P})(\mathfrak{v})(\tau)=\sum u(c_{\mu})(p_{\mu})\) .

On the other hand, let \(\nu : C \otimes_A \operatorname{Hom}_A(P, J) \longrightarrow \operatorname{Homgr}_A(\operatorname{Homgr}_A(C, P), J)\) be the homomorphism that maps the element \(c \otimes h\) , for \(c \in C_p\) , \(h \in \operatorname{Hom}_A(P, J)\) on the homomorphism \(u \mapsto (-1)^p h(u(c))\) . It is graded of degree 0; it is bijective if each module \(C_p\) is free of finite type. One verifies without difficulty that it is a morphism of complexes.

Consider the sequence of homomorphisms

\[
\begin{array}{r l} \operatorname{Tor} ^ {\Lambda} (M, \operatorname{Hom} _ {\Lambda} (P, J)) & \xrightarrow {\psi} H (C \otimes_ {A} \operatorname{Hom} _ {\Lambda} (P, J)) \xrightarrow {H (\nu)} H (\operatorname{Homgr} _ {A} (\operatorname{Homgr} _ {A} (C, P), J)) \\ & \xrightarrow {w} \operatorname{Homgr} _ {\Lambda} (H (\operatorname{Homgr} _ {A} (C, P)), J) \xrightarrow {t} \operatorname{Homgr} _ {A} (\operatorname{Ext} _ {A} (M, P), J) \end{array}
\]

where \(\psi\) is the canonical isomorphism \(\psi(\mathrm{C},\mathrm{Hom}_{\mathrm{A}}(\mathrm{P},\mathrm{J}))\) , \(w\) the canonical homomorphism \(\lambda(\mathrm{Homgr}_{\mathrm{A}}(\mathrm{C},\mathrm{P}),\mathrm{J})\) (A, X, p.~82) and \(t\) is deduced from the canonical isomorphism \(\varphi(\mathrm{C},\mathrm{P})\) . One sees as above that the composed homomorphism

\[
\rho (\mathrm{M}, \mathrm{P}): \operatorname{Tor} ^ {\mathrm{A}} (\mathrm{M}, \operatorname{Hom} _ {\mathrm{A}} (\mathrm{P}, \mathrm{J})) \longrightarrow \operatorname{Homgr} _ {\mathrm{A}} \left(\operatorname{Ext} _ {\mathrm{A}} (\mathrm{M}, \mathrm{P}), \mathrm{J}\right)
\]

is independent of the choice of the resolution C; it is \(\operatorname{End}_{\mathrm{A}}(\mathrm{J})\) -linear. Let p be an integer, \(\xi \in \operatorname{Tor}_{p}^{\mathrm{A}}(\mathrm{M}, \operatorname{Hom}_{\mathrm{A}}(\mathrm{P}, \mathrm{J}))\) , \(\lambda \in \operatorname{Ext}_{\mathrm{A}}^{p}(\mathrm{M}, \mathrm{P}), \mathrm{J}\) ; if \(\xi\) is represented by means of \(\psi(\mathrm{C}, \operatorname{Hom}_{\Lambda}(\mathrm{P}, \mathrm{J}))\) by an element \(\sum c_{\mu} \otimes u_{\mu}\) of \(\mathrm{C} \otimes \operatorname{Hom}_{\Lambda}(\mathrm{P}, \mathrm{J})\) such that \(\sum d_{\mathrm{C}}(c_{\mu}) \otimes u_{\mu} = 0\) , and by \(\lambda\) by means of \(\varphi(\mathrm{C}, \mathrm{P})\) by a homomorphism \(\ell : C_{p} \to P\) such that \(\ell \circ d_{C} = 0\) , one has \(\rho(\mathrm{M}, \mathrm{P})(\xi)(\lambda) = (-1)^{p} \sum u_{\mu} (\ell(c_{\mu}))\) .

PROPOSITION 6. - Suppose the A-module J is injective; for every A-module N, set D(N) = \(\mathrm{Hom}_{\Lambda}(N,J)\) .

\begin{enumerate}
\def\labelenumi{\alph{enumi})}
\item
  Homomorphisms \(\theta^{i}(\mathrm{M},\mathrm{P}):\operatorname{Ext}_{\mathrm{A}}^{i}(\mathrm{M},\mathrm{D}(\mathrm{P}))\longrightarrow \mathrm{D}\big(\operatorname {Tor}_{i}^{\mathrm{A}}(\mathrm{M},\mathrm{P})\big)\) are bijective.
\item
  If the ring A is Noetherian and the A-module M is finitely generated, the homomorphisms \(\rho_{i}(\mathrm{M},\mathrm{P}):\operatorname{Tor}_{i}^{\mathrm{A}}(\mathrm{M},\mathrm{D}(\mathrm{P}))\longrightarrow \mathrm{D}\bigl (\operatorname {Ext}_{\mathrm{A}}^{i}(\mathrm{M},\mathrm{P})\bigr)\) are bijective.
\item
  By construction, the homomorphism \(\theta(\mathrm{M},\mathrm{P})\) is bijective as soon as \(\lambda(\mathrm{C}\otimes_{\mathrm{A}}\mathrm{P},\mathrm{J})\) is bijective, which is the case when J is injective (A, X, p.~85, cor. 2).
\item
  Choose the resolution C so that each module \(C_{p}\) is free of finite type (A, X, p.~53, prop. 6). Then the homomorphism v is bijective, and the same is true of \(\lambda(\text{Homgr}_{\text{A}}(\text{C}, \text{P}), \text{J})\), since J is injective; therefore \(\rho(\text{M}, \text{P})\) is bijective.
\end{enumerate}

Remarks.--- 1) For every homomorphism \(f: \mathrm{N} \to \mathrm{N}'\) of A-modules, denote \(\mathrm{D}(f): \mathrm{D}(\mathrm{N}') \to \mathrm{D}(\mathrm{N})\) the homomorphism \(\mathrm{Hom}(f, \mathbf{1}_{\mathrm{J}})\) . Let \(u: \mathrm{M} \to \mathrm{M}'\) and \(v: \mathrm{P} \to \mathrm{P}'\) homomorphisms of A-modules. Choose projective resolutions (C, p) of M and (C', p') of M', and a morphism of complexes \(\tilde{u}: \mathrm{C} \to \mathrm{C}'\) such that \(p' \circ \tilde{u} = u \circ p\) (A, X, p.~49, prop. 3). The diagram

\[
\begin{array}{c c c} \operatorname{Homgr} _ {\mathrm{A}} (\mathrm{C} ^ {\prime} \otimes_ {\mathrm{A}} \mathrm{P} ^ {\prime}, \mathrm{J}) & \xrightarrow {\mu^ {\prime}} & \operatorname{Homgr} _ {\mathrm{A}} (\mathrm{C} ^ {\prime}, \operatorname{Hom} _ {\mathrm{A}} (\mathrm{P} ^ {\prime}, \mathrm{J})) \\ \operatorname{Hom} (\tilde {u} \otimes v, 1 _ {\mathrm{J}}) \Bigg \downarrow & & \Bigg \downarrow \operatorname{Hom} (\tilde {u}, \operatorname{Hom} (v, 1)) \\ \operatorname{Homgr} _ {\mathrm{A}} (\mathrm{C} \otimes_ {\mathrm{A}} \mathrm{P}, \mathrm{J}) & \xrightarrow {\mu} & \operatorname{Homgr} _ {\mathrm{A}} (\mathrm{C}, \operatorname{Hom} _ {\mathrm{A}} (\mathrm{P}, \mathrm{J})) \end{array}
\]

where \(\mu\) and \(\mu'\) are the canonical homomorphisms, is commutative; one then deduces from A. X, p.~103, prop. 2 a commutative diagram

\[
\begin{array}{c c c} \operatorname{Ext} _ {\Lambda} ^ {i} (\mathrm{M} ^ {\prime}, \mathrm{D} (\mathrm{P} ^ {\prime})) & \xrightarrow {\theta^ {i} (\mathrm{M} ^ {\prime} , \mathrm{P} ^ {\prime})} & \mathrm{D} (\operatorname{Tor} _ {i} ^ {\mathrm{A}} (\mathrm{M} ^ {\prime}, \mathrm{P} ^ {\prime})) \\ \operatorname{Ext} ^ {i} (u, \mathrm{D} (v)) \Bigg \downarrow & & \Bigg \downarrow \mathrm{D} (\operatorname{Tor} _ {i} (u, v)) \\ \operatorname{Ext} _ {\Lambda} ^ {i} (\mathrm{M}, \mathrm{D} (\mathrm{P})) & \xrightarrow {\theta^ {i} (\mathrm{M} , \mathrm{P})} & \mathrm{D} (\operatorname{Tor} _ {i} ^ {\Lambda} (\mathrm{M}, \mathrm{P})) \end{array}
\]

Let \(w: \mathrm{P}'' \to \mathrm{P}\) a homomorphism of A-modules; in an analogous manner we obtain a commutative diagram

\[
\begin{array}{c c c} \operatorname{Tor} _ {i} ^ {\mathrm{A}} (\mathrm{M}, \mathrm{D} (\mathrm{P})) & \xrightarrow {\rho_ {i} (\mathrm{M} , \mathrm{P})} & \mathrm{D} (\operatorname{Ext} _ {\mathrm{A}} ^ {i} (\mathrm{M}, \mathrm{P})) \\ \operatorname{Tor} _ {i} (u, \mathrm{D} (w)) \Bigg \downarrow & & \Bigg \downarrow \mathrm{D} (\operatorname{Ext} ^ {i} (u, w)) \\ \operatorname{Tor} _ {i} ^ {\mathrm{A}} (\mathrm{M} ^ {\prime}, \mathrm{D} (\mathrm{P} ^ {\prime \prime})) & \xrightarrow {\rho_ {i} (\mathrm{M} ^ {\prime} , \mathrm{P} ^ {\prime \prime})} & \mathrm{D} (\operatorname{Ext} _ {\mathrm{A}} ^ {i} (\mathrm{M} ^ {\prime}, \mathrm{P} ^ {\prime \prime})) \end{array} .
\]

\begin{enumerate}
\def\labelenumi{\arabic{enumi})}
\setcounter{enumi}{1}
\tightlist
\item
  Let
\end{enumerate}

\[
0 \to \mathrm{M} ^ {\prime} \xrightarrow {j} \mathrm{M} \xrightarrow {q} \mathrm{M} ^ {\prime \prime} \to 0\tag{ℰ}
\]

an exact sequence of A-modules. The homomorphism \(\mathrm{L}(q):\mathrm{L}(\mathrm{M})\to\mathrm{L}(\mathrm{M}^{\prime\prime})\) induced on the canonical free resolutions is surjective, and the complex \(\operatorname{Ker}\mathrm{L}(q)\) defines a projective resolution of \(M^{\prime}\) Applying prop. 3 of A, X, p.~104 to the exact sequence \(0\to\operatorname{Ker}\mathrm{L}(q)\to\mathrm{L}(\mathrm{M})\to\mathrm{L}(\mathrm{M}^{\prime\prime})\to0\) , we obtain commutative diagrams

\[
\begin{array}{c c c} \operatorname{Ext} _ {\mathrm{A}} ^ {i} (\mathrm{M} ^ {\prime}, \mathrm{D} (\mathrm{P})) & \xrightarrow {\theta^ {i} (\mathrm{M} ^ {\prime} , \mathrm{P})} & \mathrm{D} (\operatorname{Tor} _ {i} ^ {\mathrm{A}} (\mathrm{M} ^ {\prime}, \mathrm{P})) \\ \delta^ {i} (\mathcal {E}, \mathrm{D} (\mathrm{P})) \Bigg \downarrow & & \Bigg \downarrow (- 1) ^ {i + 1} \mathrm{D} (\partial_ {i + 1} (\mathcal {E}, \mathrm{P})) \\ \operatorname{Ext} _ {\mathrm{A}} ^ {i + 1} (\mathrm{M} ^ {\prime \prime}, \mathrm{D} (\mathrm{P})) & \xrightarrow {\theta^ {i + 1} (\mathrm{M} ^ {\prime \prime} , \mathrm{P})} & \mathrm{D} (\operatorname{Tor} _ {i + 1} ^ {\mathrm{A}} (\mathrm{M} ^ {\prime \prime}, \mathrm{P})) \end{array}
\]

\[
\begin{array}{c c c} \operatorname{Tor} _ {i + 1} ^ {\mathrm{A}} (\mathrm{M} ^ {\prime \prime}, \mathrm{D} (\mathrm{P})) & \xrightarrow {\rho_ {i + 1} (\mathrm{M} ^ {\prime \prime} , \mathrm{P})} & \mathrm{D} (\operatorname{Ext} _ {\mathrm{A}} ^ {i + 1} (\mathrm{M} ^ {\prime \prime}, \mathrm{P})) \\ \partial_ {i + 1} (\mathcal {E}, \mathrm{D} (\mathrm{P})) \Bigg \downarrow & & \Bigg \downarrow (- 1) ^ {i + 1} \mathrm{D} (\delta^ {i} (\mathcal {E}, \mathrm{P})) \\ \operatorname{Tor} _ {i} ^ {\mathrm{A}} (\mathrm{M} ^ {\prime}, \mathrm{D} (\mathrm{P})) & \xrightarrow {\rho_ {i} (\mathrm{M} ^ {\prime} , \mathrm{P})} & \mathrm{D} (\operatorname{Ext} _ {\mathrm{A}} ^ {i} (\mathrm{M} ^ {\prime}, \mathrm{P})) \end{array}
\]

Let

\[
0 \to \mathrm{P} ^ {\prime} \to \mathrm{P} \to \mathrm{P} ^ {\prime \prime} \to 0\tag{F}
\]

an exact sequence of A-modules; since the A-module J is injective, one deduces an exact sequence

\[
0 \to \mathrm{D} (\mathrm{P} ^ {\prime \prime}) \to \mathrm{D} (\mathrm{P}) \to \mathrm{D} (\mathrm{P} ^ {\prime}) \to 0.\tag{\(\left(\mathcal{D}(\mathcal{F})\right)\)}
\]

Applying A, X, p.~104, prop. 3 and p.~106, prop. 4 to the exact sequences (F) and (D(F)), one obtains in an analogous manner commutative diagrams.

\[
\begin{array}{c c c} \operatorname{Ext} _ {\Lambda} ^ {i} (\mathrm{M}, \mathrm{D} (\mathrm{P} ^ {\prime})) & \xrightarrow {\theta^ {i} (\mathrm{M} , \mathrm{P} ^ {\prime})} & \mathrm{D} (\operatorname{Tor} _ {i} ^ {\mathrm{A}} (\mathrm{M}, \mathrm{P} ^ {\prime})) \\ \delta^ {i} (\mathrm{M}, \mathrm{D} (\mathcal {F})) \Bigg \downarrow & & \Bigg \downarrow (- 1) ^ {i + 1} \mathrm{D} (\partial_ {i + 1} (\mathrm{M}, \mathcal {S})) \\ \operatorname{Ext} _ {\mathrm{A}} ^ {i + 1} (\mathrm{M}, \mathrm{D} (\mathrm{P} ^ {\prime \prime})) & \xrightarrow {\theta^ {i + 1} (\mathrm{M} , \mathrm{P} ^ {\prime \prime})} & \mathrm{D} (\operatorname{Tor} _ {i + 1} ^ {\mathrm{A}} (\mathrm{M}, \mathrm{P} ^ {\prime \prime})) \end{array}
\]

\[
\begin{array}{c c c} \operatorname{Tor} _ {i + 1} ^ {\mathrm{A}} (\mathrm{M}, \mathrm{D} (\mathrm{P} ^ {\prime})) & \xrightarrow {\rho_ {i + 1} (\mathrm{M} , \mathrm{P} ^ {\prime})} & D (\operatorname{Ext} _ {\Lambda} ^ {i + 1} (\mathrm{M}, \mathrm{P} ^ {\prime})) \\ \partial_ {i + 1} (\mathrm{M}, \mathrm{D} (\mathcal {F})) \Bigg \downarrow & & \Bigg \downarrow (\sim 1) ^ {i} \mathrm{D} (\delta^ {i} (\mathrm{M}, \mathcal {F})) \\ \operatorname{Tor} _ {i} ^ {\mathrm{A}} (\mathrm{M}, \mathrm{P} ^ {\prime \prime}) & \xrightarrow {\rho_ {i} (\mathrm{M} , \mathrm{P} ^ {\prime \prime})} & D (\operatorname{Ext} _ {\mathrm{A}} ^ {i} (\mathrm{M}, \mathrm{P} ^ {\prime \prime})) \end{array}
\]

\subsection{Dualizing modules}

\subsection{Dualizing modules}

DEFINITION 1.--- Let A be a Noetherian ring. A Λ-module Ω is said to be dualizing if it is finitely generated and if, for every maximal ideal m of Λ, the A/m-vector space Ext \(_{A}^{i}\) (A/m, Ω) is zero for i ≠ ht(m) and of dimension 1 for i = ht(m).

For every maximal ideal m of A and every integer i, the A/m-vector space \(\operatorname{Ext}_{\mathrm{A}}^{i}(\mathrm{A}/\mathfrak{m},\Omega)\) is canonically isomorphic to \(\operatorname{Ext}_{\mathrm{A}_{\mathfrak{m}}}^{i}(\mathrm{A}/\mathfrak{m},\Omega_{\mathfrak{m}})\) By consequence, for a finitely generated A-module \(\Omega\) be dualizing, it is necessary and sufficient that the \(A_{m}\) -module \(\Omega_{m}\) so that it is dualizing for every maximal ideal m of A.

Examples.--- 1) If the ring A is local and Artinian, the dualizing A-modules are the finitely generated injective A-modules. \(\Omega\) such that \(\mathrm{Hom}_{\mathrm{A}}(\kappa_{\mathrm{A}},\Omega)\) Let \( E \) be of dimension 1 ( \(\S\) 3, no. 3, prop. 6), that is, the Matlis A-modules ( \(\S\) 8, no. 3).

\begin{enumerate}
\def\labelenumi{\arabic{enumi})}
\setcounter{enumi}{1}
\tightlist
\item
  For a Noetherian ring A to be Gorenstein, it is necessary and sufficient that the A-module A be dualizing (§ 3, no. 7, prop. 11). In particular, the A-module A is dualizing when A is regular.
\end{enumerate}

Remarks. -- 1) Let A be a Noetherian local ring and \(\Omega\) a finitely generated A-module. The residue field \(\kappa_{\widehat{A}}\) is identified with \(\kappa_{A}\) , and the \(\widehat{A}\) -module \(\widehat{\Omega}\) with \(\widehat{A} \otimes_{A} \Omega\) (III, § 3, no. 4, th. 3). It then follows from A, X, p.~111, prop. 10 that the \(\kappa_{A}\) -vector space \(\mathrm{Ext}_{\mathbf{A}}^{i}(\kappa_{\mathbf{A}}, \Omega)\) is canonically isomorphic to \(\mathrm{Ext}_{\widehat{A}}^{i}(\kappa_{\widehat{A}}, \widehat{\Omega})\) . Consequently, for the \(A\) -module \(\Omega\) to be dualizing, it is necessary and sufficient that the \(\widehat{A}\) -module \(\widehat{\Omega}\) be dualizing.

\begin{enumerate}
\def\labelenumi{\arabic{enumi})}
\setcounter{enumi}{1}
\tightlist
\item
  Let \(\Omega\) a dualizing A-module; for every rank-1 projective A-module L, the A-module \(\Omega\otimes_{\mathrm{A}}\mathrm{L}\) is dualizing (A, X, p.~108, prop. 7, b)). We shall see below ( \(n^{\circ}\) 4, prop. 6) that every dualizing A-module is isomorphic to a module of this form.
\end{enumerate}

PROPOSITION 1.--- Let A be a Noetherian ring and Ω a dualizing A-module.

\begin{enumerate}
\def\labelenumi{\alph{enumi})}
\tightlist
\item
  Then A is a Macaulay ring, and the A-module Ω is Macaulayan.
\end{enumerate}

\[
\text { b) } O n a \mathrm{di} _ {\mathrm{A}} (\Omega) = \dim (\Omega) = \dim (\mathrm{A}).
\]

Suppose first that the ring A is local, and denote \(d\) its dimension. Prop. 6 of § 3, no. 3 implies \(\mathrm{di}_{\Lambda}(\Omega) = d\) , hence \(\mathrm{prof}(\mathrm{A}) = d\) by Prop. 9 of § 3, no. 6, so that A is a Macaulay ring. Moreover, we have \(\mathrm{prof}(\Omega) = d\) by definition of depth; since we have \(\mathrm{prof}(\Omega) \leqslant \dim(\Omega) \leqslant d\) from which we deduce the proposition in this case.

In the general case, the \(A_{m}\) -module \(\Omega_{m}\) is dualizing for every maximal ideal m of A, hence \(A_{m}\) is a Macaulay ring and \(\Omega_{m}\) a \(\Lambda_{m}\) a). Moreover, we have \(\mathrm{di}_{\mathrm{A}_{\mathfrak{m}}(\Omega_{\mathfrak{m}})} = \dim(\Omega_{\mathfrak{m}}) = \dim(\Lambda_{\mathfrak{m}})\) for every maximal ideal m, whence b) by passing to the supremum (§ 3, no. 2, prop. 3).

PROPOSITION 2.--- Let A be a Noetherian ring, Ω a dualizing A-module. For every prime ideal p of A, the A- \(_{p}\) -module Ω \(_{p}\) is dualizing.

Consider a saturated chain \(\mathfrak{p} \subset \mathfrak{p}_1 \subset \ldots \subset \mathfrak{p}_r\) of prime ideals of A such that the ideal \(\mathfrak{p}_r\) so is maximal. Proceeding by induction on \(r\) one may assume that the \(\mathrm{A}_{\mathfrak{p}_1}\) -module \(\Omega_{\mathfrak{p}_1}\) is dualizing. Replacing \(\Lambda\) by \(\mathrm{A}_{\mathfrak{p}_1}\) and \(\mathfrak{p}\) by \(\mathfrak{p}\mathrm{A}_{\mathfrak{p}_1}\) , we reduce to the case where the ring A is local and the chain \(\mathfrak{p} \subset \mathfrak{m}_{\Lambda}\) is saturated.

Let us then set \(d = \dim(A) = \mathrm{ht}(\mathfrak{m}_{\mathbf{A}})\) . One has \(\dim(A_{\mathfrak{p}}) = \mathrm{ht}(\mathfrak{p}) = d - 1\) since A is a Macaulay ring (§ 2, no. 2, cor. of prop. 2). For every integer \(i\) , the \(A_{\mathfrak{p}}\) -module \(\operatorname{Ext}_{A_{\mathfrak{p}}}^{i}(\kappa(\mathfrak{p}), \Omega_{\mathfrak{p}})\) is isomorphic to \(\operatorname{Ext}_{\Lambda}^{i}(A/\mathfrak{p}, \Omega)_{\mathfrak{p}}\) By § 3, no. 2, prop. 2, it therefore suffices to prove that the \(A/\mathfrak{p}\) -module \(\operatorname{Ext}_{A}^{i}(A/\mathfrak{p}, \Omega)\) is zero for \(i \neq d - 1\) and of rank one for \(i = d - 1\) .

Let \(x\) be an element of \(\mathfrak{m}_{\mathrm{A}} - \mathfrak{p}\) , and by \(\overline{x}\) its class in \(A / \mathfrak{p}\) . Consider the exact sequence of A-modules

\[
0 \rightarrow \mathrm{A} / \mathfrak {p} \xrightarrow {\overline {{{x}}}} \mathrm{A} / \mathfrak {p} \longrightarrow \mathrm{A} / (\mathfrak {p} + x \mathrm{A}) \rightarrow 0.
\]

The A-module \(\mathrm{A} / (\mathfrak{p} + x\mathrm{A})\) is of finite length since its support is reduced to \(\mathfrak{m}_{\mathrm{A}}\) ; since the A-module \(\Omega\) is dualizing, we have \(\operatorname{Ext}_{\mathrm{A}}^{i}(\mathrm{A} / (\mathfrak{p} + x\mathrm{A}), \Omega) = 0\) for \(i \neq d\) (§ 8, no. 5, example 3). One then deduces from the exact sequence of extension modules associated with the above sequence and with \(\Omega\) that the homothety of ratio \(x\) in the A-module \(\operatorname{Ext}_{\mathrm{A}}^{i}(\mathrm{A} / \mathfrak{p}, \Omega)\) is surjective for \(i \neq d - 1\) , which implies that this module is zero (Nakayama's lemma). In particular \(\operatorname{Ext}_{\mathrm{A}}^{d}(\mathrm{A} / \mathfrak{p}, \Omega)\) is zero, and one obtains an exact sequence

\[
0 \rightarrow \operatorname{Ext} _ {\Lambda} ^ {d - 1} (\mathrm{A} / \mathfrak {p}, \Omega) \xrightarrow {x} \operatorname{Ext} _ {\Lambda} ^ {d - 1} (\mathrm{A} / \mathfrak {p}, \Omega) \longrightarrow \operatorname{Ext} _ {\Lambda} ^ {d} (\mathrm{A} / (\mathfrak {p} + x \Lambda), \Omega) \rightarrow 0.
\]

We have \(\operatorname{long}_{\mathrm{A}}(\operatorname{Ext}_{\mathrm{A}}^{i}(\mathrm{A}/(\mathfrak{p}+x\mathrm{A}),\Omega))=\operatorname{long}_{\mathrm{A}}(\mathrm{A}/(\mathfrak{p}+x\mathrm{A}))\) (loc. cit.); the proposition then follows from the following lemma, applied to the ring B = A/p and to the B-module \(M=\operatorname{Ext}_{\Lambda}^{d-1}(\mathrm{A}/\mathfrak{p},\Omega)\) :

Lemma 1. Let B be a local, integral, noetherian ring of dimension 1, and let M be a finitely generated torsion-free B-module. Suppose that one has long \(_{B}\) (M/xM) = long \(_{B}\) (B/xB) for every nonzero element x of B. Then the B-module M has rank 1.

Indeed, let \(r\) the rank of M; there exists a submodule L of M free of rank \(r\) such that M/L is a torsion module (VII, § 4, no. 1, cor. of prop. 1), hence of finite length (VII, § 2, no. 5, lemma 1). The annihilator of M/L is not reduced to 0, and therefore contains a nonzero element \(x\) of \(\mathfrak{m}_{\mathbb{R}}\) . Consider the diagram

commutative

\[
\begin{array}{c c c c c c c} 0 \to & \mathrm{L} & \longrightarrow & \mathrm{M} & \longrightarrow & \mathrm{M} / \mathrm{L} & \to 0 \\ & \Biggl \downarrow x _ {\mathrm{L}} & & \Biggl \downarrow x _ {\mathrm{M}} & & \Biggl \downarrow 0 \\ 0 \to & \mathrm{L} & \longrightarrow & \mathrm{M} & \longrightarrow & \mathrm{M} / \mathrm{L} & \to 0. \end{array}
\]

By the snake lemma (A, X, p.~4, prop. 2), one deduces an exact sequence

\[
0 \rightarrow \mathrm{M} / \mathrm{L} \longrightarrow \mathrm{L} / x \mathrm{L} \longrightarrow \mathrm{M} / x \mathrm{M} \longrightarrow \mathrm{M} / \mathrm{L} \rightarrow 0,
\]

whence long(M/xM) = long(L/xL). Since long(M/xM) = long(B/xB) by hypothesis and long(L/xL) = r long(B/xB), one deduces r = 1.

COROLLARY 1. For every multiplicative subset S of A, the \(\mathrm{S}^{-1}\mathrm{A}\) -module \(\mathrm{S}^{-1}\Omega\) is dualizing.

COROLLARY 2.--- The support of Ω is equal to Spec(A).

Indeed a dualizing module over a local ring is nonzero by definition.

COROLLARY 3. Let M be a finitely generated A-module, and let i be an integer. The A-module \(\operatorname{Ext}_{\mathrm{A}}^{i}(\mathrm{M},\Omega)\) is finitely generated, and its support has codimension \(\geqslant i\) in \(\operatorname{Spec}(\mathrm{A})\) .

The first assertion follows from A, X, p.~108, cor. Let \(\mathfrak{p}\) a prime ideal of the support of \(\operatorname{Ext}_{\mathrm{A}}^{i}(\mathrm{M},\Omega)\) . One has \(\operatorname{Ext}_{\mathrm{A}}^{i}(\mathrm{M},\Omega)_{\mathfrak{p}} \neq 0\) , hence \(\operatorname{Ext}_{\mathrm{A}_{\mathfrak{p}}}^{i}(\mathrm{M}_{\mathfrak{p}},\Omega_{\mathfrak{p}}) \neq 0\) ( \(\S 3, \mathrm{n}^{\circ}2\) , prop. 2), which implies \(\mathrm{di}_{\mathrm{A}_{\mathfrak{p}}}(\Omega_{\mathfrak{p}}) \geqslant i\) . Since \(\Omega_{\mathfrak{p}}\) is a \(\mathrm{A}_{\mathfrak{p}}\) By the dualizing module (prop. 2), we have \(\mathrm{di}_{\mathrm{A}_{\mathfrak{p}}}(\Omega_{\mathfrak{p}}) = \dim(\mathrm{A}_{\mathfrak{p}})\) (prop. 1), whence the corollary.

PROPOSITION 3.--- Let A be a Noetherian local ring, Ω a dualizing A-module, and M a finitely generated A-module.

\[
\text { a) } O n a \operatorname{Ext} _ {\mathrm{A}} ^ {i} (\mathrm{M}, \Omega) = 0 \text { for } i <   \dim (\mathrm{A}) \dots \dim_ {\mathrm{A}} (\mathrm{M}).
\]

\begin{enumerate}
\def\labelenumi{\alph{enumi})}
\setcounter{enumi}{1}
\tightlist
\item
  Set \(c = \dim(\mathrm{A}) - \dim_{\mathrm{A}}(\mathrm{M})\) If M is nonzero, the A-module \(\operatorname{Ext}_{\mathrm{A}}^{c}(\mathrm{M}, \Omega)\) is not zero.
\end{enumerate}

\[
\mathrm{c)} O n a \operatorname{Ext} _ {\mathrm{A}} ^ {i} (\mathrm{M}, \Omega) = 0 p o u r i > \dim (\mathrm{A}) - \operatorname{prof} _ {\mathrm{A}} (\mathrm{M}).
\]

Suppose M is nonzero and denote by F its support. By prop. 9 of § 1, no. 5, the conjunction of assertions a) and b) is equivalent to \(\operatorname{prof}_{\mathbb{F}}(\Omega)=c\) Now since \(\Omega\) is Macaulaycn and that its support is equal to \(\operatorname{Spec}(\mathbb{A})\) (prop. 1 and cor. 2 of prop. 2), we have

\[
\operatorname{prof} _ {\mathrm{F}} (\Omega) = \operatorname{codim} (\mathrm{F}, \operatorname{Spec} (\Lambda)) = c
\]

(§ 2, no. 1, cor. of prop. 1 and no. 2, cor. of prop. 2).

Let us prove c) by induction on the depth of M. If \(\mathrm{prof}_{\Lambda}(\mathbf{M}) = 0\) , we indeed have \(\mathrm{Ext}_{\Lambda}^{i}(\mathbf{M},\Omega) = 0\) for \(i > \dim (\mathbf{A})\) , since \(\mathrm{di}_{\Lambda}(\Omega) = \dim (\mathbf{A})\) (prop. 1). Suppose \(\mathrm{prof}_{\Lambda}(\mathbf{M}) > 0\) then there exists an element \(x\) of \(\mathfrak{m}_{\mathbf{A}}\) such that the homothety of ratio \(x\) injective in M. We have \(\mathrm{prof}_{\Lambda}(\mathbf{M} / x\mathbf{M}) = \mathrm{prof}_{\Lambda}(\mathbf{M}) - 1\) ( \(\S 1, \mathrm{n}^{\circ}4\) prop. 7).

Consider the exact sequence of extension modules

\[
\operatorname{Ext} _ {\mathrm{A}} ^ {i} (\mathrm{M}, \Omega) \xrightarrow {x} \operatorname{Ext} _ {\mathrm{A}} ^ {i} (\mathrm{M}, \Omega) \longrightarrow \operatorname{Ext} _ {\Lambda} ^ {i + 1} (\mathrm{M} / x \mathrm{M}, \Omega)
\]

associated with the exact sequence

\[
0 \to \mathrm{M} \xrightarrow {x} \mathrm{M} \longrightarrow \mathrm{M} / x \mathrm{M} \to 0.
\]

For \(i > \dim(\mathrm{A}) - \operatorname{prof}_{\mathrm{A}}(\mathrm{M})\) , the A-module \(\operatorname{Ext}_{\mathrm{A}}^{i+1}(\mathrm{M}/x\mathrm{M},\Omega)\) is zero by the induction hypothesis, hence the homothety of ratio \(x\) is surjective in \(\operatorname{Ext}_{\mathrm{A}}^{i}(\mathrm{M},\Omega)\) , which implies that this A-module is zero (Nakayama's lemma). This proves c).

COROLLARY.--- If M is Macaulay, one has Ext \(^{i}\) A(M, Ω) = 0 for i ≠ c; the A-module Ext \(^{c}\) A(M, Ω) is Macaulay, and its support is equal to that of M.

The first assertion follows from prop. 3, a) and c). Let \(\mathfrak{p} \in \operatorname{Supp}(M)\) ; by prop. 1 of § 2, no. 1, applied to M and to A, we have

\[
\dim (A _ {p}) - \dim_ {A _ {p}} (M _ {p}) = \dim (A) - \dim_ {A} (M) = c;
\]

since the \(A_{p}\) -module \(\Omega_{p}\) is dualizing (prop. 2), it follows from prop. 3, b) that the \(A_{p}\) -module \(\operatorname{Ext}_{A_{p}}^{c}(M_{p},\Omega_{p})\) is not zero. Consequently, the support of \(\operatorname{Ext}_{A}^{c}(M,\Omega)\) is equal to that of M.

Let us finally prove, by induction on \(\dim(M)\) that the A-module \(\operatorname{Ext}_{A}^{c}(M,\Omega)\) is Macaulay. The assertion is satisfied when \(\dim(M)=0\) since every module of finite length is Macaulay. Suppose \(\dim(M)>0\) and let us choose an element \(x\) of \(\mathfrak{m}_{\mathrm{A}}\) such that the homothety \(x_{\mathrm{M}}\) so that it is injective. The A-module M/ \(x\mathrm{M}\) is Macaulay \(\S 2\) (no. 3, prop. 4), of dimension \(\dim(M)-1\) . In view of the above, the exact sequence of extension modules associated with the exact sequence \(0\to\mathrm{M}\xrightarrow{x}\mathrm{M}\longrightarrow\mathrm{M}/x\mathrm{M}\to0\) reduces to

\[
0 \to \operatorname{Ext} _ {\mathrm{A}} ^ {c} (\mathrm{M}, \Omega) \xrightarrow {x} \operatorname{Ext} _ {\mathrm{A}} ^ {c} (\mathrm{M}, \Omega) \longrightarrow \operatorname{Ext} _ {\mathrm{A}} ^ {c + 1} (\mathrm{M} / x \mathrm{M}, \Omega) \to 0;
\]

Proposition 4 of §2, No. 3 and the induction hypothesis then imply that Ext \(_{A}^{c}\) (M, Ω) is Macaulay, whence the corollary.

\subsection{Quotient by a regular sequence}

PROPOSITION 4. --- Let A be a Noetherian ring, J an ideal of A generated by an A-regular sequence x, and Ω a finitely generated A-module.

\begin{enumerate}
\def\labelenumi{\alph{enumi})}
\item
  If the A-module \(\Omega\) is dualizing, the sequence \(\mathbf{x}\) is \(\Omega\) -regular and the A/J-module \(\Omega/\mathrm{J}\Omega\) is dualizing;
\item
  If the A/J-module \(\Omega/\mathrm{J}\Omega\) is dualizing, that J is contained in the radical of \(\Lambda\) and that the sequence x is \(\Omega\) -regular, the A-module \(\Omega\) is dualizing.
\end{enumerate}

Reasoning by induction on the length of the sequence x, we reduce to the case where it is reduced to a single element x. Suppose that the A-module \(\Omega\) let it be dualizing. For every maximal ideal m of A containing x, we have \(\dim(\mathrm{A}_{\mathfrak{m}}/x\mathrm{A}_{\mathfrak{m}})=\dim(\mathrm{A}_{\mathfrak{m}})-1\) (VIII, § 3, no. 1, cor. 2), and consequently \(\mathrm{Hom}_{\mathrm{A}_{\mathfrak{m}}}\left(\mathrm{A}_{\mathfrak{m}} / x\mathrm{A}_{\mathfrak{m}}, \Omega_{\mathfrak{m}}\right) = 0\) (no. 1, prop. 3, a)). This implies \(\mathrm{Hom}_{\mathrm{A}}(\mathrm{A} / x\mathrm{A}, \Omega) = 0\) so that the homothety \(x_{\Omega}\) is injective. One may therefore assume, in order to prove the proposition, that the homothety \(x_{\Omega}\) is injective.

Let \(\overline{\mathsf{A}}\) the ring A/xA; let \(\mathfrak{m}\) a maximal ideal of A containing \(x\) , and let \(\overline{\mathfrak{m}}\) its image in \(\overline{\mathsf{A}}\) . The A-module A/ \(\mathfrak{m}\) is annihilated by \(x\) and is identified with \(\overline{\mathsf{A}}/\overline{\mathfrak{m}}\) ; thus for every integer \(i \geqslant 1\) we have an isomorphism \(\operatorname{Ext}_{\mathsf{A}}^{i}(\mathsf{A}/\mathfrak{m}, \Omega) \longrightarrow \operatorname{Ext}_{\overline{\mathsf{A}}}^{i-1}(\overline{\mathsf{A}}/\overline{\mathfrak{m}}, \Omega/x\Omega)\) ( \(\S\) 3, no. 4, prop. 7). We have

\[
\mathrm{ht} (\overline {{\mathfrak {m}}}) = \dim (\overline {{\mathrm{A}}} _ {\overline {{\mathfrak {m}}}}) = \dim (\mathrm{A} _ {\mathfrak {m}} / x \mathrm{A} _ {\mathfrak {m}}) = \dim (\mathrm{A} _ {\mathfrak {m}}) - 1 = \mathrm{ht} (\mathfrak {m}) - 1
\]

(VIII, § 3, no. 1, cor. 2, a)). Now the maximal ideals of \(\overline{\mathrm{A}}\) are the ideals \(\overline{\mathfrak{m}}\) , where \(\mathfrak{m}\) is a maximal ideal of A containing \(x\) if moreover \(x\) belongs to the radical of A, every maximal ideal of A contains \(x\) The proposition follows from this.

COROLLARY 1. Let A be an integral noetherian ring. Every dualizing A-module is torsion-free and of rank 1.

Let \(\Omega\) a dualizing A-module; it is torsion-free by prop. 4. Let K be the field of fractions of A; the K-vector space \(K \otimes_{\Lambda} \Omega\) is dualizing \(n^{\circ}\) 1, prop. 2), hence of dimension 1.

COROLLARY 2. Let A be a local Macaulay ring, Ω a finitely generated A-module, and x a maximal secant sequence of elements of m \(_{A}\) , generating an ideal J. The following conditions are equivalent:

\begin{enumerate}
\def\labelenumi{(\roman{enumi})}
\item
  the A-module Ω is dualizing;
\item
  the A-module \(\Omega\) is Macaulay of dimension equal to \(\dim(A)\) , and by \(\Omega/J\Omega\) is an indecomposable injective module over the local artinian ring \(A/J\) ;
\item
  the sequence x is Ω-regular and Ω/JΩ is an indecomposable injective module over the local artinian ring A/J ;
\item
  the sequence \(\mathbf{x}\) is \(\Omega\) -regular, one has \(\mathrm{long}_{\mathrm{A}}(\Omega/\mathrm{J}\Omega) = \mathrm{long}_{\mathrm{A}}(\Lambda/\mathrm{J})\) and the \(\kappa_{\mathrm{A}}\) -vector space \(\mathrm{Hom}_{\Lambda}(\kappa_{\Lambda}, \Omega/\mathrm{J}\Omega)\) has dimension 1.
\end{enumerate}

(i)⇒(ii): if Ω is dualizing, it is Macaulay and of dimension dim(A) (n° 1, prop. 1). The sequence x is A-regular since A is a Macaulay ring; by prop. 4, the A/J-module Ω/JΩ is dualizing, hence is a Matlis A/J-module (n° 1, example 1).

(ii)⇒(iii): under hypothesis (ii), one has \(\dim(\Omega)=\dim(A)\) and \(\dim(\Omega/J\Omega)=\dim(A/J)=0\) , so that the sequence x is secant for \(\Omega\) , hence \(\Omega\) -regular (§ 2, n° 3, th. 1).

(iii)⇒(i): under the hypotheses of (iii), the Λ/J-module Ω/JΩ is a Matlis A-module, hence is dualizing (n° 1, example 1); by prop. 4 the A-module Ω is dualizing.

\begin{enumerate}
\def\labelenumi{(\roman{enumi})}
\setcounter{enumi}{2}
\tightlist
\item
  \(\Leftrightarrow\) (iv): this follows from the remark of § 8, n° 3.
\end{enumerate}

\subsection{Change of rings}

PROPOSITION 5.--- Let \(\rho: A \to B\) a homomorphism of noetherian rings, making B a flat A-module. One assumes that for every maximal ideal \(\mathfrak{n}\) of B, the ring \(\kappa(\rho^{-1}(\mathfrak{n})) \otimes_{A} B\) is a Gorenstein ring. Let \(\Omega\) be a dualizing A-module; the B-module \(\Omega_{(B)}\) is dualizing.

Let \(\mathfrak{n}\) a maximal ideal of B, and \(\mathfrak{p}\) its inverse image in A. The \(A_{\mathfrak{p}}\) -module \(B_{\mathfrak{n}}\) is flat, the \(A_{\mathfrak{p}}\) -module \(\Omega_{\mathfrak{p}}\) is dualizing, \(\Omega_{(B)} \otimes_B B_{\mathfrak{n}}\) is identified with \(\Omega_{\mathfrak{p}} \otimes_{A_{\mathfrak{p}}} B_{\mathfrak{n}}\) and \(\kappa_{\Lambda_{\mathfrak{p}}} \otimes_{\Lambda_{\mathfrak{p}}} B_{\mathfrak{n}}\) , which is identified with a ring of fractions of \(\kappa(\mathfrak{p}) \otimes_A B\) , is a Gorenstein ring. It therefore suffices to prove the proposition when \(\rho\) is a local homomorphism of local rings, which we will assume from now on.

Let us first treat the case where the rings A and B are Artinian. Set \(C = B / \mathfrak{m}_{\Lambda}B\) Since B is flat over A, the B-module \(\mathrm{Hom}_{\mathrm{B}}(\mathrm{C},\Omega_{(\mathrm{B})})\) is isomorphic to \(\mathrm{Hom}_{\Lambda}(\kappa_{\mathrm{A}},\Omega)\otimes_{\mathrm{A}}\mathrm{B}\) (I, § 2, no. 10, prop. 11), hence to \(\mathrm{Hom}_{\mathrm{A}}(\kappa_{\mathrm{A}},\Omega)\otimes_{\kappa_{\mathrm{A}}}\mathrm{C}\) From this we deduce a sequence of isomorphisms

\[
\begin{array}{c} \operatorname{Hom} _ {\mathrm{B}} (\kappa_ {\mathrm{B}}, \Omega_ {(\mathrm{B})}) \longrightarrow \operatorname{Hom} _ {\mathrm{C}} (\kappa_ {\mathrm{C}}, \operatorname{Hom} _ {\mathrm{B}} (\mathrm{C}, \Omega_ {(\mathrm{B})})) \longrightarrow \operatorname{Hom} _ {\mathrm{C}} (\kappa_ {\mathrm{C}}, \operatorname{Hom} _ {\mathrm{A}} (\kappa_ {\mathrm{A}}, \Omega) \otimes_ {\kappa_ {\mathrm{A}}} \mathrm{C}) \\ \qquad \qquad \qquad \qquad \qquad \qquad \qquad \qquad \qquad \qquad \longrightarrow \operatorname{Hom} _ {\mathrm{A}} (\kappa_ {\mathrm{A}}, \Omega) \otimes_ {\kappa_ {\mathrm{A}}} \operatorname{Hom} _ {\mathrm{C}} (\kappa_ {\mathrm{C}}, \mathrm{C}). \end{array}
\]

The \(\kappa_{\mathrm{A}}\) -vector space \(\mathrm{Hom}_{\mathrm{A}}(\kappa_{\mathrm{A}}, \Omega)\) is of dimension 1 since \(\Omega\) is dualizing, and the same holds for the \(\kappa_{\mathrm{C}}\) -vector space \(\mathrm{Hom}_{\mathrm{C}}(\kappa_{\mathrm{C}}, \mathrm{C})\) since C is a Gorenstein ring; consequently, \(\kappa_{\mathrm{B}}\) -vector space \(\mathrm{Hom}_{\mathrm{B}}(\kappa_{\mathrm{B}}, \Omega_{(\mathrm{B})})\) has dimension 1.

Let M be a B-module of finite length; let us prove by induction on \(\operatorname{long}_{\mathrm{B}}(\mathrm{M})\) that we have \(\operatorname{long}_{\mathrm{B}}(\operatorname{Hom}_{\mathrm{B}}(\mathrm{M},\Omega_{(\mathrm{B})})) \leqslant \operatorname{long}_{\mathrm{B}}(\mathrm{M})\) The assertion is clear if M = 0, and it follows from what precedes if \(M = \kappa_{B}\) . Suppose \(\operatorname{long}_{\mathrm{B}}(\mathrm{M}) \geqslant 2\) There exists an exact sequence of B-modules

\[
0 \rightarrow \mathrm{M} ^ {\prime} \rightarrow \mathrm{M} \rightarrow \kappa_ {\mathrm{B}} \rightarrow 0
\]

with \(\mathrm{long}_{\mathrm{B}}(\mathrm{M}^{\prime}) < \mathrm{long}_{\mathrm{B}}(\mathrm{M})\) One deduces an exact sequence

\[
0 \to \operatorname{Hom} _ {\mathrm{B}} (\kappa_ {\mathrm{B}}, \Omega_ {(\mathrm{B})}) \to \operatorname{Hom} _ {\mathrm{B}} (\mathrm{M}, \Omega_ {(\mathrm{B})}) \to \operatorname{Hom} _ {\mathrm{B}} (\mathrm{M} ^ {\prime}, \Omega_ {(\mathrm{B})}),
\]

and one concludes by applying the induction hypothesis to M'.

Let N be the kernel of the canonical surjection from \(\kappa_{A} \otimes_{A} B\) on \(\kappa_{B}\) . Set \(m = \text{long}_{\text{B}}(\kappa_{\text{A}} \otimes_{\text{A}} \text{B})\) ; one has \(\text{long}_{\text{B}}(\text{N}) = m - 1\) . Consider the exact sequence of B-modules

\[
\begin{array}{c} 0 \to \operatorname{Hom} _ {\mathrm{B}} (\kappa_ {\mathrm{B}}, \Omega_ {(\mathrm{B})}) \longrightarrow \operatorname{Hom} _ {\mathrm{B}} (\kappa_ {\mathrm{A}} \otimes_ {\mathrm{A}} \mathrm{B}, \Omega_ {(\mathrm{B})}) \longrightarrow \operatorname{Hom} _ {\mathrm{B}} (\mathrm{N}, \Omega_ {(\mathrm{B})}) \\ \longrightarrow \operatorname{Ext} _ {\mathrm{B}} ^ {1} (\kappa_ {\mathrm{B}}, \Omega_ {(\mathrm{B})}) \longrightarrow \operatorname{Ext} _ {\mathrm{B}} ^ {1} (\kappa_ {\mathrm{A}} \otimes_ {\mathrm{A}} \mathrm{B}, \Omega_ {(\mathrm{B})}). \end{array}
\]

The B-modules \(\mathrm{Hom}_{\mathbb{B}}(\kappa_{\Lambda}\otimes_{\Lambda}\mathrm{B},\Omega_{(\mathbb{B})})\) and \(\mathrm{Ext}_{\mathbb{B}}^{1}(\kappa_{\Lambda}\otimes_{\Lambda}\mathrm{B},\Omega_{(\mathbb{B})})\) are respectively isomorphic to \(\mathrm{Hom}_{\Lambda}(\kappa_{\Lambda},\Omega)\otimes_{\Lambda}\mathrm{B}\) and \(\mathrm{Ext}_{\mathbb{A}}^{1}(\kappa_{\Lambda},\Omega)\otimes_{\Lambda}\mathrm{B}\) , that is, to \(\kappa_{\Lambda}\otimes_{\Lambda}\mathrm{B}\) and to 0. The lengths of the B-modules \(\mathrm{Hom}_{\mathbb{B}}(\kappa_{\mathbb{B}},\Omega_{(\mathbb{B})})\) and \(\mathrm{Hom}_{\mathbb{B}}(\kappa_{\mathbb{A}}\otimes_{\mathbb{A}}\mathrm{B},\Omega_{(\mathbb{B})})\) are 1 and \(m\) and that of \(\mathrm{Hom}_{\mathbb{B}}(\mathrm{N},\Omega_{(\mathbb{B})}))\) is \(\leqslant m - 1\) from which it follows that the B-module \(\mathrm{Ext}_{\mathbb{B}}^{1}(\kappa_{\mathbb{B}},\Omega_{(\mathbb{B})})\) is zero. By prop. 6 of § 3, no. 3, the B-module \(\Omega_{(\mathbb{B})}\) is injective; consequently it is a dualizing module (no. 1, example 1).

Moving to the general case, we set \(C = \kappa_A \otimes_A B\) ; by hypothesis, this is a Gorenstein ring, hence a Macaulay ring ( \(\S\) By prop. 1 of no. 1, A is a Macaulay ring, and the A-module \(\Omega\) is Macaulay. Consequently B is a Macaulay ring, and the B-module \(\Omega_{(B)}\) is Macaulay \(\S\) 2, no. 7, cor. 1 of prop. 9). Set \(r = \dim(A)\) , \(s = \dim(C)\) There exists a sequence ( \(x_1, \ldots, x_r\) ) of elements of \(\mathfrak{m}_A\) regular for the A-modules A and \(\Omega\) and a sequence ( \(y_1, \ldots, y_s\) ) of elements of \(\mathfrak{m}_B\) regular for the B-module C; denote \(x\) denote the ideal of A and \(y\) the ideal of B that they generate respectively. The sequence ( \(y_1, \ldots, y_s, \rho(x_1), \ldots, \rho(x_r)\) ) is regular for the B-modules B and \(\Omega_{(B)}\) ( \(\S\) 1, no. 6, prop. 11), and the A-module B/ \(y\) is flat (loc. cit., prop. 10). Set \(A' = A/x\) , \(B' = B/(xB + y)\) and let \(\rho': A' \to B'\) the homomorphism deduced from \(\rho\) by passing to quotients. The rings \(A'\) and \(B'\) are Artinian, the A'-module \(B'\) is flat, the ring \(\kappa_{A'} \otimes_{A'} B'\) , which is identified with C/ \(y\) is a Gorenstein ring ( \(\S\) 3, no. 7, example 2) and the A'-module \(\Omega_{(A')}\) is dualizing (no. 2, prop. 4). By the first part of the proof, the B'-module \(\Omega_{(B')}\) is dualizing. It then follows from loc. cit. that the B-module \(\Omega_{(B)}\) is dualizing.

COROLLARY.---Let A be a Noetherian ring admitting a dualizing module. \(\Omega\) B be a polynomial algebra over A in a finite number of indeterminates. The B-module \(\Omega_{(B)}\) is dualizing.

Indeed, for every prime ideal p of A, the ring \(\kappa(\mathfrak{p})\otimes_{\mathrm{A}}\mathrm{A}[\mathbf{X}]\) is identified with \(\kappa(\mathfrak{p})[\mathbf{X}]\) , which is regular, hence Gorenstein.

PROPOSITION 6. — Let \(\Lambda\) be a Noetherian local ring and \(\Omega\) a dualizing A-module. Let B be a finite A-algebra; assume that the A-module B is Macaulay. The B-module \(\mathrm{Ext}_{\mathrm{A}}^{i}(\mathrm{B},\Omega)\) is zero for \(i\neq \dim (\Lambda) - \dim (\mathrm{B})\) and dualizing for \(i = \dim (\Lambda) - \dim (\mathrm{B})\) .

We have \(\dim(B) = \dim_{A}(B) \leqslant \dim(\Lambda)\) (VIII, § 2, no. 3, th. 1 c)); set \(c = \dim(A) - \dim(B)\) . One has \(\operatorname{Ext}_{A}^{i}(B, \Omega) = 0\) for \(i \neq c\) since the A-module B is Macaulay (no. 1, cor. of prop. 3). Let us prove that the B-module \(\operatorname{Ext}_{\Lambda}^{c}(B, \Omega)\) is dualizing.

Suppose first \(\dim(B) = 0\) . The spectrum X of B is finite and consists of maximal ideals (IV, § 2, no. 5, prop. 9); the canonical map \(B \to \prod_{n \in X} B_n\) is an isomorphism (loc. cit., cor. 1). The B-module \(\Omega' = \text{Ext}_A^c(B, \Omega)\) is therefore a direct sum of the modules \(\text{Ext}_A^c(B_n, \Omega)\) ; as \(\text{Ext}_A^c(B_n, \Omega)\) has support in \(\{\mathfrak{n}\}\) , it is identified with \(\Omega'_n\) . One has \(\dim(B_n) = 0\) for every \(\mathfrak{n}\) ; to prove that the B-module \(\Omega'\) is dualizing, it therefore suffices to prove that this is so for the \(B_n\) -module \(\text{Ext}_A^c(B_n, \Omega)\) for every \(\mathfrak{n} \in X\) , which reduces us to the case where the ring B is local. In this case, according to example 6 of § 8, no. 5, the B-module \(\text{Ext}_A^c(B, \Omega)\) is isomorphic to \(\text{Hom}_A(B, I)\) , where I is a Matlis A-module; it is consequently a Matlis B-module (§ 8, no. 6, cor. of prop. 5), hence a dualizing B-module (no. 1, example 1).

Now suppose \(\dim(B) > 0\) and reason by induction on \(\dim(B)\) . One has \(\operatorname{prof}_{A}(B) = \dim_{A}(B) = \dim(B)\) , whence \(\operatorname{prof}_{A}(B) > 0\) ; on the other hand we have \(\operatorname{prof}(A) = \dim(A) > 0\) (no. 1, prop. 1), and consequently \(\operatorname{prof}_{A}(A \oplus B) > 0\) . There therefore exists an element \(x\) of \(\mathfrak{m}_{\Lambda}\) such that the homotheties \(x_{\Lambda}\) and \(x_{B}\) are injective.

Consider the exact sequence of extension modules associated with the exact sequence \(0 \to \mathrm{B} \xrightarrow{x_{\mathrm{B}}} \mathrm{B} \longrightarrow \mathrm{B}/x\mathrm{B} \to 0\) and with the A-module \(\Omega\) . The A-module B/xB is Macaulay ( \(\S 2, n^{\circ} 1\) , example 3), of dimension \(\dim(\mathrm{B}) - 1\) (VIII, \(\S 3, n^{\circ} 2\) ; we therefore have \(\operatorname{Ext}_{\mathrm{A}}^{i}(\mathrm{B}/x\mathrm{B}, \Omega) = 0\) for \(i \neq c + 1\) ( \(n^{\circ} 1\) (corollary of Prop. 3). Since we have \(\operatorname{Ext}_{\mathrm{A}}^{i}(\mathrm{B}, \Omega) = 0\) for \(i \neq c\) , one obtains an exact sequence of B-modules

\[
0 \rightarrow \operatorname{Ext} _ {\mathrm{A}} ^ {c} (\mathrm{B}, \Omega) \xrightarrow {x} \operatorname{Ext} _ {\mathrm{A}} ^ {c} (\mathrm{B}, \Omega) \longrightarrow \operatorname{Ext} _ {\mathrm{A}} ^ {c + 1} (\mathrm{B} / x \mathrm{B}, \Omega) \rightarrow 0.
\]

By the induction hypothesis, the B/xB-module \(\operatorname{Ext}_{\Lambda}^{c+1}(\mathrm{B}/x\mathrm{B},\Omega)\) is dualizing. Since the A-algebra B is finite, the image of \(m_{A}\) in B is contained in the radical of B (V, § 2, no. 1, prop. 1); by prop. 4 of no. 2, the B-module \(\operatorname{Ext}_{\mathrm{A}}^{c}(\mathrm{B},\Omega)\) is dualizing.

COROLLARY 1.--- Let A be a Noetherian ring, Ω a dualizing A-module, and B a finite A-algebra; suppose that the A-module B is Macaulay. The B-module Ext \(_{A}\) (B, Ω) is dualizing.

Let \(\Omega'\) the B-module \(\mathrm{Ext}_{\mathsf{A}}(\mathsf{B},\Omega)\) . Let \(\mathfrak{n}\) a maximal ideal of B; its inverse image in A is a maximal ideal \(\mathfrak{m}\) (V, § 2, no. 1, prop. 1). The \(\Lambda_{\mathfrak{m}}\) -algebra \(\mathrm{B}_{\mathfrak{m}} = \mathrm{A}_{\mathfrak{m}} \otimes_{\mathsf{A}} \mathrm{B}\) is finite, and it is a \(\mathrm{A}_{\mathfrak{m}}\) -module macaulayen; by the proposition, the \(\mathrm{B}_{\mathfrak{m}}\) -module \(\Omega_{\mathfrak{m}}'\) , which is identified with \(\mathrm{Ext}_{\mathrm{A}_{\mathfrak{m}}}(B_{\mathfrak{m}},\Omega_{\mathfrak{m}})\) (§ 3, no. 2, prop. 2) is dualizing. Since \(\mathrm{B}_{\mathfrak{n}}\) is a ring of fractions of \(\mathrm{B}_{\mathfrak{m}}\) , the \(\mathrm{B}_{\mathfrak{n}}\) -module \(\Omega_{\mathfrak{n}}'\) is dualizing, whence the corollary.

Remark. Keep the hypotheses of Cor. 1 and suppose further that the canonical homomorphism \(\rho : A \to B\) be injective. We then have \(\dim(\Lambda_{\mathfrak{m}}) = \dim(B_{\mathfrak{m}})\) for every maximal ideal \(\mathfrak{m}\) of \(A\) By prop. 6 and cor. 1, \(\operatorname{Ext}_{A}^{i}(B, \Omega)\) is zero for \(i \neq 0\) and the B-module \(\operatorname{Hom}_{A}(B, \Omega)\) is dualizing.

COROLLARY 2.--- If a Noetherian ring A possesses a dualizing module, then every finitely generated A-algebra which is a Macaulay ring possesses a dualizing module.

This follows from cor. 1 and from the corollary of prop. 5.

COROLLARY 3.--- Every presentable Macaulay ring (in particular, every complete local Macaulay ring) possesses a dualizing module.

Indeed, let R be a regular ring and A a Macaulay ring quotient of R. The R-module A is Macaulay (§ 2, no. 5, example 5), and R possesses a dualizing module (no. 1, example 2); hence so does Λ by cor. 1. Moreover, it has already been observed that a complete local Noetherian ring is presentable (§ 4, no. 4, prop. 6, c)).

More generally, every Macaulay ring that is a quotient of a Gorenstein ring has a dualizing module. Conversely, one can show that a local Macaulay ring that has a dualizing module is a quotient of a local Gorenstein ring (exerc. 1).

\subsection{Structure of dualizing modules}

Lemma 2.--- Let A be a Noetherian ring, and let M and N be finitely generated A-modules, \(u: \mathrm{M} \to \mathrm{N}\) a homomorphism. Let \(x\) an element of the radical of A, such that the homothety \(x_{\mathrm{N}}\) be injective. If the homomorphism \(\overline{u}: \mathrm{M}/x\mathrm{M} \to \mathrm{N}/x\mathrm{N}\) induced by \(u\) is injective (resp. surjective, resp. bijective), then so is \(u\) .

The assertion concerning the surjectivity of \(u\) follows from Nakayama's lemma (II, § 3, no. 2, cor. 1 of prop. 4), without any hypothesis on \(x_{\mathrm{N}}\) . Consider the commutative diagram with exact rows

\[
\begin{array}{c c c c c c c} & \mathrm{M} & \xrightarrow {x _ {\mathrm{M}}} & \mathrm{M} & \longrightarrow & \mathrm{M} / x \mathrm{M} \\ & u \Bigg \downarrow & & u \Bigg \downarrow & & \overline {{u}} \Bigg \downarrow \\ 0 & \longrightarrow & \mathrm{N} & \xrightarrow {x _ {\mathrm{N}}} & \mathrm{N} & \longrightarrow & \mathrm{N} / x \mathrm{N} \end{array} ;
\]

using the snake lemma (I, § 1, n° 4, prop. 2), one deduces an exact sequence \(\operatorname{Ker} u \xrightarrow{x} \operatorname{Ker} u \longrightarrow \operatorname{Ker} \overline{u}\) If \(\overline{u}\) is injective, the homothety with ratio \(x\) is surjective in \(\operatorname{Ker} u\) , which implies \(\operatorname{Ker} u = 0\) by Nakayama's lemma.

PROPOSITION 7.--- Let A be a Noetherian ring and Ω a dualizing A-module.

\begin{enumerate}
\def\labelenumi{\alph{enumi})}
\item
  We have \(\operatorname{Ext}_{\mathrm{A}}^{i}(\Omega, \Omega) = 0\) for \(i > 0\) .
\item
  \(L'\) canonical homomorphism \(\gamma : \mathrm{A} \to \operatorname{End}_{\mathrm{A}}(\Omega)\) is bijective.
\item
  Every dualizing A-module is of the form \(\Omega\otimes_{\mathrm{A}}\mathrm{L}\) where \(\mathrm{L}\) is a projective A-module of rank 1.
\end{enumerate}

\begin{enumerate}
\def\labelenumi{\Alph{enumi})}
\tightlist
\item
  Let us first treat the case where the ring A is local. In this case, condition c) simply means that two dualizing modules are isomorphic.
\end{enumerate}

Let \(\Omega'\) a dualizing A-module. We have \(\operatorname{prof}_{\mathrm{A}}(\Omega') = \dim_{\mathrm{A}}(\Omega') = \dim(\mathrm{A})\) (no. 1, prop. 1), hence \(\operatorname{Ext}_{\Lambda}^{i}(\Omega', \Omega) = 0\) for \(i \neq 0\) (no. 1, prop. 3, c)), whence a).

Let us prove b) and c) by induction on the integer dim(A) (equal to prof(A)). If it is zero, the ring A is Artinian, \(\Omega'\) and \(\Omega\) are Matlis A-modules ( \(n^{\circ}\) 1, example 1); they are therefore isomorphic ( \(\S\) 8, \(n^{\circ}\) the canonical map A → \(\mathrm{End}_{\mathrm{A}}(\Omega)\) is bijective \(\S\) 8, \(n^{\circ}\) 2, prop. 3, c)). Suppose dim(A) \textgreater{} 0 and let \(x\) be a cancellable element of \(\mathfrak{m}_{\mathrm{A}}\) . The homothety \(x_{\Omega}\) is injective ( \(n^{\circ}\) 2, prop. 4), and one has an exact sequence

\[
0 \rightarrow \Omega \xrightarrow {x _ {\Omega}} \Omega \longrightarrow \Omega / x \Omega \rightarrow 0.
\]

Since \(\mathrm{Ext}_{\mathrm{A}}^{1}(\Omega',\Omega)\) is zero and that \(\mathrm{Hom}_{\mathrm{A}}(\Omega',\Omega /x\Omega)\) is identified with \(\mathrm{Hom}_{\mathrm{A} / x\mathrm{A}}(\Omega'/x\Omega',\Omega/x\Omega)\) , one deduces an exact sequence

\[
0 \rightarrow \operatorname{Hom} _ {\mathrm{A}} \left(\Omega^ {\prime}, \Omega\right) \xrightarrow {x} \operatorname{Hom} _ {\mathrm{A}} \left(\Omega^ {\prime}, \Omega\right) \xrightarrow {p} \operatorname{Hom} _ {\mathrm{A} / x \mathrm{A}} \left(\Omega^ {\prime} / x \Omega^ {\prime}, \Omega / x \Omega\right)\rightarrow 0,\tag{1}
\]

where \(p\) is the canonical map. By prop. 4, the A/xA-modules \(\Omega/x\Omega\) and \(\Omega'/x\Omega'\) are dualizing, hence isomorphic by the induction hypothesis. Let \(\overline{u}\) an isomorphism of \(\Omega'/x\Omega'\) on \(\Omega/x\Omega\) . Taking into account the exact sequence (1), there exists an A-homomorphism \(u: \Omega' \to \Omega\) such that \(p(u) = \overline{u}\) ; by Lemma 2, \(u\) is bijective, which proves c). By the induction hypothesis, the canonical homomorphism A/xA \(\longrightarrow\) End \(_{A/xA}(\Omega/x\Omega)\) is bijective. Given the exact sequence (1), this homomorphism is identified with the homomorphism \(\overline{\gamma}: A/xA \longrightarrow \text{End}_A(\Omega)/x \text{End}_A(\Omega)\) induced by \(\gamma\) it then follows from Lemma 2 that \(\gamma\) is bijective, whence b).

\begin{enumerate}
\def\labelenumi{\Alph{enumi})}
\setcounter{enumi}{1}
\tightlist
\item
  Let us turn to the general case. For every maximal ideal m of A and every integer i \textgreater{} 0, we have \(\operatorname{Ext}_{A_{m}}^{i}(\Omega_{\mathfrak{m}}, \Omega_{\mathfrak{m}}) = 0\) based on the foregoing, therefore \(\operatorname{Ext}_{A}^{i}(\Omega, \Omega)_{\mathfrak{m}} = 0\) (§ 3, no. 2, prop. 2), which implies \(\operatorname{Ext}_{A}^{i}(\Omega, \Omega) = 0\) (II, § 3, no. 3, cor. 2 of th. 1). Likewise, the homomorphism \(\gamma_{m}: A_{m} \to \operatorname{End}_{A}(\Omega)_{m}\) is bijective for every maximal ideal m of A, hence \(\gamma\) is bijective (loc. cit., th. 1).
\end{enumerate}

Let us finally prove c). Let \(\Omega'\) a dualizing A-module. Denote by L the A-module \(\mathrm{Hom}_{\mathrm{A}}(\Omega',\Omega)\) and by \(v:\Omega' \otimes_{\mathrm{A}} \mathrm{L} \longrightarrow \Omega\) the homomorphism such that \(v(x \otimes f) = f(x)\) for \(x \in \Omega'\) , \(f \in \mathrm{L}\) . Let \(\mathfrak{m}\) a maximal ideal of A. The \(\mathrm{A}_{\mathfrak{m}}\) -module \(\mathrm{L}_{\mathfrak{m}}\) is identified with \(\mathrm{Hom}_{\mathrm{A}_{\mathfrak{m}}}(\Omega_{\mathfrak{m}}',\Omega_{\mathfrak{m}})\) by the already treated case it is free of rank one, and every isomorphism \(h:\Omega_{\mathfrak{m}}' \to \Omega_{\mathfrak{m}}\) is a generator. When one identifies \(\mathrm{L}_{\mathfrak{m}}\) with \(\mathrm{A}_{\mathfrak{m}}\) using the generator \(h\) , the homomorphism \(v_{\mathfrak{m}}:\Omega_{\mathfrak{m}}' \otimes_{\mathrm{A}_{\mathfrak{m}}}\mathrm{L}_{\mathfrak{m}} \longrightarrow \Omega_{\mathfrak{m}}\) is identified with \(h\) therefore it is bijective. Since this holds for every maximal ideal \(\mathfrak{m}\) of A, the A-module L is projective of rank one (II, § 5, no. 3, th. 2), and the homomorphism \(v\) is bijective (II, § 3, no. 3, th. 1).

COROLLARY 1.--- For A to be a Gorenstein ring, it is necessary and sufficient that the A-module Ω be projective of rank 1.

This follows from example 2 of no. 1 and from prop. 7 c).

COROLLARY 2.--- Suppose that the ring A is presentable. The set of prime ideals p of A such that A \(_{p}\) let a Gorenstein ring be open in Spec(A).

Let p be a prime ideal of A such that \(A_{p}\) Let A be a Gorenstein ring. It is then a Macaulay ring; after possibly replacing A by \(A_{f}\) , where f is a suitable element of A - p, one reduces to the case where A is a Macaulay ring (§ 4, n° 4, prop. 7, c)). Let \(\Omega\) be a dualizing A-module (n° 3, cor. 3 of prop. 6). Then \(\Omega_{p}\) is a dualizing module over the Gorenstein ring \(A_{p}\) (n° 1, prop. 2), hence is free of rank 1 (cor. 1). Consequently there exists an element g of A - p such that the \(A_{g}\) -module \(\Omega_{g}\) is free of rank 1 (II, § 5, n° 1, cor. of prop. 2). Thus \(A_{g}\) is a Gorenstein ring (cor. 1) and the same is true of \(A_{q}\) for every prime ideal q of A not containing g (§ 3, n° 7, example 1), which proves the corollary.

\subsection{Duality of finitely generated modules}

In this section we consider a Noetherian ring A of finite dimension which possesses a dualizing module \(\Omega\) . One then has \(\mathrm{di}_{\Lambda}(\Omega) = \dim(\mathrm{A}) < +\infty\) ( \(n^{\circ}\) 1, prop. 1). Choose an injective resolution of finite length \(e: \Omega \to (\mathbf{I}, \delta)\) . For every complex C of A-modules, denote \(\mathbf{D}(\mathbf{C})\) the complex \(\operatorname{Homgr}_{\Lambda}(\mathbf{C}, \mathbf{I})\) . This applies in particular to every A-module M, considered as a complex concentrated in degree 0; we then have \(\mathbf{D}(\mathrm{M})^{i} = \mathrm{Hom}_{\mathrm{A}}(\mathrm{M},\mathrm{I}^{i})\) for every integer \(i\) . Recall that in A, X, p.~100, th. 1, we constructed a canonical isomorphism

\[
\varphi (\mathrm{M}, \mathrm{I}): \mathrm{H} (\mathbf {D} (\mathrm{M})) \longrightarrow \operatorname{Ext} _ {\mathrm{A}} (\mathrm{M}, \Omega).
\]

Examples.--- 1) The complex \(\mathbf{D}(\mathbf{A}) = \operatorname{Homgr}_{\mathbf{A}}(\mathbf{A},\mathbf{l})\) is identified with I. The map \(c:\Omega \to \mathbf{D}(\Lambda)\) is by definition a homologism.

\begin{enumerate}
\def\labelenumi{\arabic{enumi})}
\setcounter{enumi}{1}
\item
  The homomorphism \(e \in \text{Homgr}_A(\Omega, I)^0\) is an element of \(\mathbf{D}(\Omega)^0\) ; the A-linear map \(\tilde{e}: A \to \mathbf{D}(\Omega)\) such that \(\tilde{e}(1) = e\) is a homologism (no. 4, prop. 7, a) and b)).
\item
  Let S be a multiplicative subset of A. The \(S^{-1}A\) -module \(S^{-1}\Omega\) is dualizing (no. 1, cor. 1 of prop. 2); the \(S^{-1}A\) -modules \(S^{-1}I^{i}\) are injective (cor. 1 of prop. 3 of § 3, no. 2) and the morphism \(e^{\prime}:S^{-1}\Omega\to S^{-1}I\) deduced from e is an injective resolution of \(S^{-1}\Omega\) , to which one may therefore apply the preceding results. For every complex C of finite type (and in particular every finitely generated A-module M), the canonical homomorphism from \(S^{-1}\mathbf{D}(\mathbf{C})\) in \(\operatorname{Homgr}_{\mathrm{S}^{-1}\mathrm{A}}(\mathrm{S}^{-1}\mathrm{C},\mathrm{S}^{-1}\mathrm{I})=\mathbf{D}(\mathrm{S}^{-1}\mathrm{C})\) is bijective.
\item
  Let A be a Dedekind ring, K its field of fractions. The A-module A is dualizing and admits the injective resolution I of length 1 defined by the exact sequence
\end{enumerate}

\[
0 \to \mathrm{A} \stackrel {e} {\longrightarrow} \mathrm{K} \stackrel {\delta} {\longrightarrow} \mathrm{K/A} \to 0
\]

where δ is the canonical surjection. For every A-module M, the complex D(M) is the complex concentrated in degrees 0 and 1

\[
\dots \longrightarrow 0 \longrightarrow \operatorname{Hom} _ {\mathrm{A}} (\mathrm{M}, \mathrm{K}) \stackrel {d} {\longrightarrow} \operatorname{Hom} _ {\mathrm{A}} (\mathrm{M}, \mathrm{K} / \mathrm{A}) \longrightarrow 0 \longrightarrow \dots
\]

with \(d = \mathrm{Hom}_{\mathrm{A}}(1_{\mathrm{M}},\delta)\) . We have an exact sequence

\[
0 \to \operatorname{Hom} _ {\mathrm{A}} (\mathrm{M}, \mathrm{A}) \longrightarrow \mathbf {D} (\mathrm{M}) ^ {0} \stackrel {d} {\longrightarrow} \mathbf {D} (\mathrm{M}) ^ {1} \longrightarrow \operatorname{Ext} _ {\mathrm{A}} ^ {1} (\mathrm{M}, \mathrm{A}) \to 0.
\]

For every morphism of complexes \(f: \mathrm{C} \to \mathrm{C}'\) , we denote \(\mathbf{D}(f): \mathbf{D}(\mathrm{C}') \to \mathbf{D}(\mathrm{C})\) the morphism of complexes \(\operatorname{Homgr}_{\mathrm{A}}(f, 1_{\mathrm{I}})\) If \(f\) is a homologism, \(\mathbf{D}(f)\) is a homologism (A, X, p.~86, prop. 4, b)). If \(\mathrm{C}' \xrightarrow{f} \mathrm{C} \xrightarrow{g} \mathrm{C}''\) is an exact sequence of complexes, the sequence of complexes \(\mathbf{D}(\mathrm{C}''') \xrightarrow{\mathbf{D}(g)} \mathbf{D}(\mathrm{C}) \xrightarrow{\mathbf{D}(f)} \mathbf{D}(\mathrm{C}')\) is exact (A, X, p.~83, prop. 2, a)).

Let M be an A-module. To each element m of M, associate the map \(\alpha_{\mathrm{M}}(m): f \mapsto f(m)\) of D(M) into I; it is an element of \(\mathbf{D}(\mathbf{D}(\mathbf{M}))_{0} = \operatorname{Homgr}_{\mathrm{A}}(\mathbf{D}(\mathbf{M}), \mathrm{I})_{0}\) It follows from the definitions that \(\alpha_{\mathrm{M}}(m)\) is a morphism of complexes, hence an element of \(\mathrm{Z}_{0}(\mathbf{D}(\mathbf{D}(\mathbf{M})))\) .

One thus defines a morphism of complexes:

\[
\alpha_ {\mathrm{M}}: \mathrm{M} \rightarrow \mathbf {D} (\mathbf {D} (\mathrm{M})) ,
\]

whence, by passing to homology, a homomorphism of A-modules

\[
\alpha_ {\mathrm{M}}: \mathrm{M} \to \mathrm{H} _ {0} (\mathbf {D} (\mathbf {D} (\mathrm{M}))).
\]

THEOREM 1. Let M be a finitely generated A-module. Then \(\alpha_{\mathrm{M}}\) is a homomorphism: we have \(\mathrm{H}_{i}(\mathbf{D}(\mathbf{D}(\mathbf{M}))) = 0\) for \(i \neq 0\) and the homomorphism \(\alpha_{\mathrm{M}}\) is bijective.

Let us first take M = A. The map \(e : \Omega \to \mathbf{D}(A)\) is a homomorphism (example 1), hence so is the map \(\mathbf{D}(e) : \mathbf{D}(\mathbf{D}(A)) \to \mathbf{D}(\Omega)\) . The map \(\tilde{e} : A \to \mathbf{D}(\Omega)\) is a homologism (example 2), and we have \(\mathbf{D}(e) \circ \alpha_{\mathbf{A}} = \tilde{c}\) ; thus \(\alpha_{A}\) is a homomorphism, which proves the theorem in this case. It follows that \(\alpha_{M}\) is a homomorphism when the A-module M is finitely generated free.

Let us turn to the general case; we shall prove by induction on the integer n the following assertion:

\((A_{n})\) for every finitely generated A-module M, the homomorphism \(H_{i}(\alpha_{M})\) is bijective for \(i \leqslant n\) .

This also means that \(\mathrm{H}_{i}(\mathbf{D}(\mathbf{D}(\mathbf{M})))\) is zero for \(i \neq 0\) and \(i \leqslant n\) , and that \(\alpha_{M}\) is bijective if \(n \geqslant 0\) . Observe that \((\mathrm{A}_{n})\) holds for n \textless{} -d, where d is the length of the complex 1: indeed the A-module \(\mathbf{D}(\mathbf{D}(\mathbf{M}))_{i}\) is equal to \(\bigoplus_{p}\mathrm{Hom}_{\mathrm{A}}(\mathrm{Hom}_{\mathrm{A}}(\mathrm{M},\mathrm{I}^{p}),\mathrm{I}^{p-i})\) , hence is zero for i \textless{} -d and i \textgreater{} d.

Let us prove the implication \((\mathrm{A}_{n}) \Rightarrow (\mathrm{A}_{n+1})\) . Let M be a finitely generated A-module. There exists a finitely generated free A-module L and an exact sequence \(0 \to N \xrightarrow{u} L \xrightarrow{v} M \to 0\) . The sequence \(0 \to \mathbf{D}(M) \xrightarrow{\mathbf{D}(v)} \mathbf{D}(L) \xrightarrow{\mathbf{D}(u)} \mathbf{D}(N) \to 0\) is exact; likewise, if we set \(u' = \mathbf{D}(\mathbf{D}(u))\) and \(v' = \mathbf{D}(\mathbf{D}(v))\) , the sequence \(0 \to \mathbf{D}(\mathbf{D}(N)) \xrightarrow{u'} \mathbf{D}(\mathbf{D}(L)) \xrightarrow{v'} \mathbf{D}(\mathbf{D}(M)) \to 0\) is exact.

Since \(H_{i}(\mathbf{D}(\mathbf{D}(\mathrm{L})))\) is zero for \(i \neq 0\) we have isomorphisms

\[
\mathrm{H} _ {i} (\mathbf {D} (\mathbf {D} (\mathrm{M}))) \longrightarrow \mathrm{H} _ {i - 1} (\mathbf {D} (\mathbf {D} (\mathrm{N}))) \quad \text { pour } i \neq 0, 1;
\]

this entails the implication \((\mathrm{A}_{n})\Rightarrow(\mathrm{A}_{n+1})\) for \(n\neq-1\) and \(n\neq0\) . Consider the commutative diagram with exact rows

\[
\begin{array}{c c c c c c c} 0 \to & N & \xrightarrow {u} & L & \xrightarrow {v} & M & \to 0 \\ & \alpha_ {N} \Bigg \downarrow & & \alpha_ {L} \Bigg \downarrow & & \alpha_ {M} \Bigg \downarrow \\ 0 \to H _ {1} (\mathbf {D} (\mathbf {D} (M))) & \xrightarrow {} H _ {0} (\mathbf {D} (\mathbf {D} (N))) & \xrightarrow {H _ {0} (u ^ {\prime})} & H _ {0} (\mathbf {D} (\mathbf {D} (L))) & \xrightarrow {H _ {0} (v ^ {\prime})} & H _ {0} (\mathbf {D} (\mathbf {D} (M))) \longrightarrow H _ {- 1} (\mathbf {D} (\mathbf {D} (N)) \end{array}
\]

where \(\alpha_{\mathrm{L}}\) is bijective. If \((\mathrm{A}_0)\) is satisfied, the homomorphism \(\alpha_{\mathrm{N}}\) is also bijective, hence \(\mathrm{H}_0(u')\) is injective and one obtains \(\mathrm{H}_1(\mathbf{D}(\mathbf{D}(\mathbf{M}))) = 0\) , whence \((\mathrm{A}_1)\) If \((\mathrm{A}_{-1})\)

is satisfied, \(\mathrm{H}_{-1}(\mathbf{D}(\mathbf{D}(\mathrm{N})))\) is zero, therefore \(\mathrm{H}_{0}(v^{\prime})\) is surjective, which implies that \(\alpha_{M}\) is surjective. Since this holds for every finitely generated A-module M, \(\alpha_{N}\) is also surjective; by I, § 1, no. 4, cor. 2 of prop. 2, \(\alpha_{M}\) is bijective, so that \((\mathrm{A}_{0})\) is satisfied.

Thus \((\mathrm{A}_{n})\) is true for all n, which proves the theorem.

Let M be a finitely generated A-module; set \(c = \dim(A) - \dim_{A}(M)\) . Let \(\mathbf{D}'(M)\) the subcomplex of \(\mathbf{D}(M)\) equal to \(\bigoplus_{i\in\mathcal{O}}\mathbf{D}(M)^{i}\bigoplus Z^{c}(\mathbf{D}(M))\) , and by

\[
j: \mathbf {D} ^ {\prime} (\mathrm{M}) \rightarrow \mathbf {D} (\mathrm{M})
\]

the canonical injection. From the canonical surjection we deduce \(\mathrm{Z}^{c}(\mathbf{D}(\mathrm{M}))\to\mathrm{H}^{c}(\mathbf{D}(\mathrm{M}))\) and of the isomorphism \(\varphi(\mathrm{M},\mathrm{I})\) a morphism of complexes

\[
p: \mathbf {D} ^ {\prime} (\mathrm{M}) (- c) \rightarrow \operatorname{Ext} _ {\Lambda} ^ {c} (\mathrm{M}, \Omega);
\]

as \(\mathrm{H}^{i}(\mathbf{D}(\mathrm{M}))\) is zero for i \textless{} c (no. 1, prop. 3 a)), \((\mathbf{D}^{\prime}(\mathrm{M})(-c), p)\) is a left resolution of \(\operatorname{Ext}_{\mathrm{A}}^{c}(\mathrm{M}, \Omega)\) . By A, X, p.~100, th. 1, we have a canonical isomorphism

\[
\varphi^ {0} \left(\mathbf {D} ^ {\prime} (\mathrm{M}) (- c), \mathrm{I}\right): \mathrm{H} _ {0} \left(\mathbf {D} \left(\mathbf {D} ^ {\prime} (\mathrm{M})\right)\right) \longrightarrow \operatorname{Ext} _ {\mathrm{A}} ^ {c} \left(\operatorname{Ext} _ {\mathrm{A}} ^ {c} (\mathrm{M}, \Omega), \Omega\right);
\]

composing it with the homomorphism \(\mathrm{H}_{0}(\mathbf{D}(j)) : \mathrm{H}_{0}(\mathbf{D}(\mathbf{D}(\mathrm{M}))) \to \mathrm{H}_{0}(\mathbf{D}(\mathbf{D}'(\mathrm{M})))\) we therefore obtain a homomorphism

\[
\mathrm{H} _ {0} (\mathbf {D} (\mathbf {D} (\mathrm{M}))) \longrightarrow \operatorname{Ext} _ {\mathrm{A}} ^ {c} \left(\operatorname{Ext} _ {\mathrm{A}} ^ {c} (\mathrm{M}, \Omega), \Omega\right),
\]

whence finally, by composition with \(\alpha_{\mathrm{M}}\) , a canonical homomorphism

\[
\beta_ {\mathrm{M}}: \mathrm{M} \longrightarrow \operatorname{Ext} _ {\mathrm{A}} ^ {c} (\operatorname{Ext} _ {\mathrm{A}} ^ {c} (\mathrm{M}, \Omega), \Omega).
\]

COROLLARY.--- If the A-module M is Macaulay, the homomorphism \(\beta_{\mathrm{M}}\) is bijective.

If M is Macaulay, the \(\Lambda\) -module \(\mathrm{H}^i(\mathbf{D}(\mathrm{M}))\) is zero for \(i \neq c\) ( \(n^\circ 1\) (cor. of prop. 3), so that the canonical injection \(j: \mathbf{D}'(\mathrm{M}) \to \mathbf{D}(\mathrm{M})\) is a homologism; consequently the morphism of complexes \(\mathbf{D}(j): \mathbf{D}(\mathbf{D}(\mathrm{M})) \to \mathbf{D}(\mathbf{D}'(\mathrm{M}))\) is a homologism ( \(\Lambda\) , X, p.~86, prop. 4). Thus \(\mathrm{H}_0(\mathbf{D}(j))\) is bijective; on the other hand, \(\alpha_{\mathrm{M}}\) is bijective by th. 1, whence the corollary.

When the A-module M is of finite length, the A-module \(\operatorname{Ext}_{\Lambda}^{c}(M,\Omega)\) is identified with the Matlis dual of M (cf.~§ 8, n° 5, example 3 and th. 3), and one recovers prop. 4 of § 8, n° 4.

\subsection{Example: the case of dimension 1}

In this section, one considers an integral, noetherian ring A of dimension 1, admitting a dualizing module \(\Omega\) . We denote by K the field of fractions of A, and by V the K-vector space \(K \otimes_{A} \Omega\) .

The canonical homomorphism \(\Omega\to V\) is injective, and the K-vector space \(V\) is of dimension 1 ( \(n^{\circ}\) 2, cor. 1 of prop. 4); let us identify \(\Omega\) to a sub-A-module of \(V\) .

PROPOSITION 8. The A-module V/Ω is a Matlis module.

Consider the exact sequence

\[
0 \to \Omega \to \mathrm{V} \to \mathrm{V} / \Omega \to 0;
\]

the A-module V is injective (A, X, p.~18, example 1), and we have \(\mathrm{di}_{\mathrm{A}}(\Omega) = 1\) ( \(n^{\circ} 1\) , prop. 1). From this we deduce on the one hand that \(V/\Omega\) is injective ( \(\S 3\) , \(n^{\circ} 1\) , prop. 1), on the other hand, that for every maximal ideal \(\mathfrak{m}\) of A, the \(\Lambda/\mathfrak{m}\) -vector space \(\mathrm{Hom}_{\mathrm{A}}(\mathrm{A}/\mathfrak{m}, \mathrm{V}/\Omega)\) is isomorphic to \(\mathrm{Ext}_{\mathrm{A}}^{1}(\mathrm{A}/\mathfrak{m}, \Omega)\) , hence of dimension 1. Since \(V/\Omega\) is a torsion module, its associated prime ideals are maximal; this proves the proposition ( \(\S 8\) , \(n^{\circ} 4\) ).

Let M be an A-module; in accordance with loc. cit., we denote by D(M) the A-module \(\mathrm{Hom}_{\mathrm{A}}(\mathrm{M},\mathrm{V}/\Omega)\) . One can apply the constructions of no. 5 by taking for I the complex

\[
\dots 0 \longrightarrow \mathrm{V} \xrightarrow {p} \mathrm{V} / \Omega \longrightarrow 0 \dots
\]

where V is placed in degree 0, and p denotes the canonical surjection. The complex D(M) is

\[
\dots 0 \longrightarrow \operatorname{Hom} _ {\mathrm{A}} (\mathrm{M}, \mathrm{V}) \xrightarrow {p _ {\mathrm{M}}} \mathrm{D} (\mathrm{M}) \longrightarrow 0 \dots ,
\]

with \(p_{\mathbf{M}} = \mathrm{Hom}(\mathbf{1}_{\mathbf{M}}, p)\) We have a canonical isomorphism of graded A-modules \(\varphi(\mathbf{M}, \mathbf{I}) : \mathrm{H}(\mathbf{D}(\mathbf{M})) \longrightarrow \mathrm{Ext}_{\mathbf{A}}(\mathbf{M}, \Omega)\) (A, X, p.~100, th. 1).

When M is a torsion module, the module \(\mathbf{D}(\mathbf{M})^{0} = \operatorname{Hom}_{\mathbf{A}}(\mathbf{M}, \mathbf{V})\) is zero, and \(\varphi(\mathbf{M}, \mathbf{I})\) is an isomorphism from D(M) to \(\operatorname{Ext}_{\Lambda}^{1}(\mathbf{M}, \Omega)\) ; the morphism \(\alpha_{M}: M \to \mathbf{D}(\mathbf{D}(\mathbf{M}))\) is none other than the canonical homomorphism of A-modules

\[
\alpha_ {\mathrm{M}}: \mathrm{M} \longrightarrow \mathrm{D} (\mathrm{D(M)})
\]

defined in § 8, no. 3, which is an isomorphism when M is of finite type (that is, of finite length). In this case we recover the situation of loc. cit.

Let us return to the general case, and suppose the A-module M is of finite type. Then D(M) is a torsion module, so the A-module \(\mathbf{D}(\mathbf{D}(\mathbf{M}))^{-1} = \mathrm{Hom}_{\mathrm{A}}(\mathrm{D}(\mathrm{M}), \mathrm{V})\) is zero. On the other hand, the A-module \(\mathrm{Hom}_{\mathrm{A}}(\mathbf{D}(\mathbf{M})^{0}, \mathrm{V}) = \mathrm{Hom}_{\mathrm{A}}(\mathrm{Hom}_{\mathrm{A}}(\mathrm{M}, \mathrm{V}), \mathrm{V})\) is naturally identified with \(K \otimes_{A} M\) , so that the homomorphism

\[
\operatorname{Hom} (1, p): \operatorname{Hom} _ {\mathrm{A}} (\operatorname{Hom} _ {\mathrm{A}} (\mathrm{M}, \mathrm{V}), \mathrm{V}) \longrightarrow \mathrm{D} (\operatorname{Hom} _ {\mathrm{A}} (\mathrm{M}, \mathrm{V}))
\]

is identified with the map

\[
j: \mathrm{K} \otimes_ {\mathrm{A}} \mathrm{M} \longrightarrow \mathrm{D} (\operatorname{Hom} _ {\mathrm{A}} (\mathrm{M}, \mathrm{V}))
\]

such that \(j(\lambda \otimes m)(f) = p(\lambda f(m))\) for \(\lambda \in \mathbf{K}\) , \(m \in \mathbf{M}\) , \(f \in \mathrm{Hom}_{\mathrm{A}}(\mathrm{M},\mathrm{V})\) . Th. 1 of no. 5 therefore translates into the exactness of the sequence

\[
0 \longrightarrow \mathrm{M} \xrightarrow {(i , \alpha_ {\mathrm{M}})} (\mathrm{K} \otimes_ {\mathrm{A}} \mathrm{M}) \oplus \mathrm{D} (\mathrm{D} (\mathrm{M})) \xrightarrow {(j , - \mathrm{D} (p _ {\mathrm{M}}))} \mathrm{D} (\operatorname{Hom} _ {\Lambda} (\mathrm{M}, \mathrm{V})) \longrightarrow 0,
\]

where \(i\) denotes the canonical map from M into \(\mathrm{K} \otimes_{\mathrm{A}} \mathrm{M}\) The kernel of \(i\) is identified with the torsion submodule \(\mathrm{T(M)}\) of M, and its cokernel at \((\mathrm{K/A}) \otimes_{\mathrm{A}} \mathrm{M}\) . Consider the commutative diagram with exact rows

\[
\begin{array}{c c c c c c c} 0 \to \mathrm{T(M)} & \longrightarrow & \mathrm{M} & \stackrel {{\imath}} {{\longrightarrow}} & \mathrm {K\otimes_ {A} M} & \longrightarrow & (\mathrm{K/A}) \otimes_ {\mathrm{A}} \mathrm{M} \longrightarrow 0 \\ & & ^ {\alpha_ {\mathrm{M}}} \Bigg \downarrow & & \Bigg \downarrow_ {j} \\ 0 \to \mathrm{D(Ext} _ {\mathrm{A}} ^ {1} (\mathrm{M}, \Omega)) & \longrightarrow & \mathrm{D(D(M))} & \stackrel {{\mathrm{D} (p _ {\mathrm{M}})}} {{\longrightarrow}} & \mathrm{D(Hom} _ {\mathrm{A}} (\mathrm{M}, \mathrm{V})) & \longrightarrow & \mathrm{D(Hom} _ {\mathrm{A}} (\mathrm{M}, \Omega)) \longrightarrow 0 \end{array}
\]

where the second line is obtained by Matlis duality from the exact sequence

\[
0 \rightarrow \operatorname{Hom} _ {\mathrm{A}} (\mathrm{M}, \Omega) \longrightarrow \operatorname{Hom} _ {\mathrm{A}} (\mathrm{M}, \mathrm{V}) \xrightarrow {p _ {\mathrm{M}}} \operatorname{Hom} _ {\mathrm{A}} (\mathrm{M}, \mathrm{V} / \Omega) \longrightarrow \operatorname{Ext} _ {\Lambda} ^ {1} (\mathrm{M}, \Omega) \rightarrow 0.
\]

Theorem 1 then means that the homomorphisms of A-modules

\[
\gamma^ {0} (\mathrm{M}): \mathrm{T} (\mathrm{M}) \longrightarrow \mathrm{D} \left(\operatorname{Ext} _ {\mathrm{A}} ^ {1} (\mathrm{M}, \boldsymbol {\Omega})\right) e t \quad \gamma^ {1} (\mathrm{M}): (\mathrm{K} / \mathrm{A}) \otimes_ {\Lambda} \mathrm{M} \longrightarrow \mathrm{D} \left(\operatorname{Hom} _ {\mathrm{A}} (\mathrm{M}, \boldsymbol {\Omega})\right)
\]

inferred from \(\alpha_{\mathrm{M}}\) and \(j\) respectively, are bijective. Since the A-module T(M) has finite length, the A-module \(\operatorname{Ext}_{\mathrm{A}}^{1}(\mathrm{M},\Omega)\) is of finite length and is identified with the Matlis dual D(T(M)), and we have \([\operatorname{Ext}_{\mathrm{A}}^{1}(\mathrm{M},\Omega)] = [\operatorname{T}(\mathrm{M})]\) in the group Z0(A) and \(\operatorname{long}_{\mathrm{A}}(\operatorname{Ext}_{\mathrm{A}}^{1}(\mathrm{M},\Omega)) = \operatorname{long}_{\mathrm{A}}(\operatorname{T}(\mathrm{M}))\) ( \(\S 8\) (no. 4, prop. 4). On the other hand, when one takes M = A, one obtains a canonical isomorphism \(\gamma^{1}(\mathrm{A}):K / A\to D(\Omega)\) .

Let B be a subring of K containing A, finite over \(\Lambda\) . For every maximal ideal \(\mathfrak{m}\) of A, we have \(\mathrm{prof}_{\mathrm{A}_{\mathfrak{m}}}(B_{\mathfrak{m}}) = \dim_{A_{\mathfrak{m}}}(B_{\mathfrak{m}}) = 1\) ( \(\S 1\) , \(n^{\circ} 1\) , Remark 2 and VIII, \(\S 2\) , \(n^{\circ} 3\) , Th. 1), so that B is a Macaulay A-module. Consequently, the B-module \(\Omega_{\mathrm{B}} = \mathrm{Hom}_{\mathrm{A}}(B, \Omega)\) is dualizing \(n^{\circ} 3\) , remark). The canonical map from \(\Omega_{\mathrm{B}} = \mathrm{Hom}_{\mathrm{A}}(B, \Omega)\) in \(\Omega = \mathrm{Hom}_{\mathrm{A}}(A, \Omega)\) is injective; its image consists of the elements \(\omega\) of \(\Omega\) such that the sub-B-module B \(\omega\) of V be contained in \(\Omega\) Thus \(\Omega_{\mathrm{B}}\) is identified with the largest sub-B-module of \(\Omega\) . The A-module B/ \(\Lambda\) is of finite length; the exact sequence

\[
0 \to \Omega_ {\mathrm{B}} \to \Omega \to \operatorname{Ext} _ {\mathrm{A}} ^ {1} (\mathrm{B/A}, \Omega) \to 0
\]

allows one to identify \(\Omega/\Omega_{\mathrm{B}}\) with \(\operatorname{Ext}_{\mathrm{A}}^{1}(\mathrm{B}/\mathrm{A},\Omega)\) , hence by what precedes with \(\mathrm{D}(\mathrm{B}/\mathrm{A})\) . In particular, we have \([\mathrm{B}/\mathrm{A}] = [\Omega/\Omega_{\mathrm{B}}]\) in \(Z_{0}(\mathrm{A})\) and \(\operatorname{Ann}_{\mathrm{A}}(\mathrm{B}/\mathrm{A}) = \operatorname{Ann}_{\mathrm{A}}(\Omega/\Omega_{\mathrm{B}})\) ( \(\S 8\) , \(\mathfrak{n}^{\circ}4\) , prop. 4).

The ideal \(\mathfrak{c} = \mathrm{Ann}_{\mathrm{A}}(\mathrm{B / A})\) is the transporter A : B, that is (VII, § 1, no. 1) the set of elements \(x\) of K such that \(xB \subset \Lambda\) . It is a (nonzero) ideal of

A and of B; in fact it is the largest ideal of B contained in A. Since \(\Omega_{\mathrm{B}} : \Omega \subset \Omega : \Omega = \mathrm{A} (\mathrm{n}^{\circ} 4, \text{prop. } 7, \mathrm{b}))\) , we have Ann \(_{\mathrm{A}}(\Omega/\Omega_{\mathrm{B}}) = \Omega_{\mathrm{B}} : \Omega\) , whence finally

\[
\mathfrak {c} = \operatorname{Ann} _ {\mathrm{A}} (\mathrm{B/A}) = \operatorname{Ann} _ {\mathrm{A}} (\Omega / \Omega_ {\mathrm{B}}) = \Omega_ {\mathrm{B}}: \Omega .
\]

Since \(\Omega_{\mathrm{B}}\) is a B-module, the relation \(x\Omega \subset \Omega_{\mathrm{B}}\) is equivalent to \(xB\Omega \subset \Omega_{\mathrm{B}}\) , so that we also have \(\mathfrak{c} = \Omega_{\mathrm{B}}: \mathrm{B}\Omega\) .

We shall now specialize the preceding discussion to the case where B is the integral closure of A; the hypothesis that B is a finitely generated A-module is satisfied when the ring A is Japanese (IX, § 4, no. 1, def. 1), which is the case when it is local and complete (loc. cit., no. 2, th. 2), or when it is essentially of finite type over a field (loc. cit., no. 1, remark 2 and example). The ring B is then a Dedekind ring (VII, § 2, no. 2, th. 1), and the torsion-free B-modules \(\Omega_{\mathrm{B}}\) , B \(\Omega\) and c are projective of rank 1 (VII, § 4, no. 10, prop. 22). The relation \(c = \Omega_{\mathrm{B}} : \mathrm{B}\Omega\) then means that the linear map \(c \otimes_{\mathrm{B}} \mathrm{B}\Omega \to \Omega_{\mathrm{B}}\) deduced from the action of K on V is an isomorphism (II, § 5, no. 6, prop. 11). In particular, we have \(\Omega_{\mathrm{B}} = c(\mathrm{B}\Omega) = c\Omega\) .

PROPOSITION 9. Let B be the integral closure of A, and \(\mathfrak{c} = \mathrm{A}:\mathrm{B}\) . Suppose that B is a finitely generated A-module. We have the inequality \([\mathrm{B} / \mathrm{c}] \leqslant 2[\mathrm{B} / \mathrm{A}]\) in \(\mathrm{Z}_0(\mathrm{A})\) . For equality to hold, it is necessary and sufficient that A be a Gorenstein ring.

We have \([B/c] = [B/A] + [A/c]\) , so that the inequality under consideration is equivalent to \([A/c] \leqslant [B/A]\) .

\begin{enumerate}
\def\labelenumi{\Alph{enumi})}
\item
  For every maximal ideal m of A, the integral closure of \(A_{m}\) is \(B_{m}\) (V, § 1, no. 5, cor. 1), and we have \(c_{m} = A_{m} : B_{m}\) . Moreover, by definition we have \([B/c] = \sum_{m} \text{long}_{A_{m}}(B_{m}/c_{m})[m]\) and \([B/A] = \sum_{m} \text{long}_{A_{m}}(B_{m}/A_{m})[m]\) . This reduces us to proving the proposition when the ring A is local, which we shall henceforth assume. In this case the ordered group \(Z_{0}(A)\) is canonically identified with Z, so that the class of a module of finite length is its length. The ring B is semi-local and the B-module BΩ is free of rank 1 (II, § 5, no. 3, prop. 5).
\item
  If A is a Gorenstein ring, the A-module \(\Omega\) is free of rank 1 (no. 4, cor. 1 of prop. 7); let us choose a generator \(\omega\) of \(\Omega\) . One has \(\Omega = A\omega\) and \(\Omega_{\mathrm{B}} = c\Omega = c\omega\) and consequently \(\text{long}(A/c) = \text{long}(\Omega/\Omega_{\mathrm{B}}) = \text{long}(B/A)\) .
\item
  Suppose the residue field \(\kappa_{\mathrm{A}}\) is infinite. For every maximal ideal \(\mathfrak{n}\) of B, denote by L(n) the B/n-vector space (of dimension 1) BΩ/nBΩ, and by pr \(_{\mathfrak{n}}\) the canonical projection of \(\bigoplus_{\mathfrak{n}}\) L(n) onto L(n). Let \(\varphi : \Omega \longrightarrow \bigoplus_{\mathfrak{n}}\) L(n) be the restriction to \(\Omega\) of the canonical homomorphism BΩ \(\longrightarrow \bigoplus_{\mathfrak{n}}\) L(n). The image of \(\varphi\) is a sub- \(\kappa_{\mathrm{A}}\) -vector space of \(\bigoplus_{\mathfrak{n}}\) L(n); it is not contained in Ker pr \(_{\mathfrak{n}}\) , for otherwise one would have \(\Omega \subset \mathfrak{n}\mathrm{B}\Omega\) and consequently BΩ \(\subset \mathfrak{n}\mathrm{B}\Omega\) , which is contradictory. Thus the image of \(\varphi\) is not contained in the union of the Ker pr \(_{\mathfrak{n}}\) (A, V, p.~40, lemma 1); there therefore exists an element \(\omega\) of \(\Omega\) whose image in BΩ/nBΩ is nonzero for every n, which implies that \(\omega\) generates the B-module BΩ (II, § 3, no. 3, prop. 11).
\end{enumerate}

Let \(a \in \mathbf{A}\) ; if \(a\omega\) belongs to \(\Omega_{\mathrm{B}}\) , one has \(a\mathrm{B}\omega \subset \Omega_{\mathrm{B}}\) , hence \(a\Omega \subset \Omega_{\mathrm{B}}\) , which implies \(a \in \mathfrak{c}\) . The map \(a \mapsto a\omega\) therefore induces an injection of \(\mathrm{A}/\mathfrak{c}\) in \(\Omega/\Omega_{\mathrm{B}}\) ; consequently one has \(\mathrm{long}(\mathrm{A}/\mathfrak{c}) \leqslant \mathrm{long}(\Omega/\Omega_{\mathrm{B}}) = \mathrm{long}(\mathrm{B}/\mathrm{A})\) .

If long(A/c) = long(B/A), then Aω + ΩB = Ω. We may assume that the ideal c is contained in mA (otherwise A equals B, hence is a Gorenstein ring). Since ΩB = cΩ is contained in mAΩ, it follows from Nakayama's lemma that ω generates Ω. Thus the A-module Ω is cyclic, hence free of rank 1, which means that A is a Gorenstein ring (no. 4, cor. 1 of prop. 7).

\begin{enumerate}
\def\labelenumi{\Alph{enumi})}
\setcounter{enumi}{3}
\tightlist
\item
  Let us treat the general case. Denote by A' the ring A{]}X{[}, that is (IX, App., no. 2) the local ring of the polynomial ring Λ{[}X{]} in the prime ideal m\_AΛ{[}X{]} ; it is a flat, integral A-algebra of dimension 1, whose residue field κ\_A' is identified with κ\_A(X) and the field of fractions to K(X) (loc. cit.). By the corollary of Prop. 5 of No. 3, the A'-module A' ⊗\_A Ω is dualizing. Set B' = A' ⊗\_A B; this is the integral closure of Λ' in K(X) (V, § 1, no. 3, prop. 13 and no. 5, prop. 16). The transporter c' = A' : B' is equal to cA' (I, § 2, no. 10, formula (11)). For every A-module M of finite length, one has long\_A'(A' ⊗\_A M) = long\_A(M): indeed, since the A-algebra A' is flat, it suffices to prove this relation when M is simple, that is, isomorphic to κ\_A; but in this case A' ⊗\_A κ\_A is identified with κ\_A′, whence our assertion. We therefore have
\end{enumerate}

\[
\operatorname{long} _ {\mathrm{A}} (\mathrm{B} / \mathfrak {c}) = \operatorname{long} _ {\mathrm{A} ^ {\prime}} \left(\mathrm{B} ^ {\prime} / \mathfrak {c} ^ {\prime}\right) \quad \text { and } \quad \operatorname{long} _ {\mathrm{A}} (\mathrm{B} / \mathrm{A}) = \operatorname{long} _ {\mathrm{A} ^ {\prime}} \left(\mathrm{B} ^ {\prime} / \mathrm{A} ^ {\prime}\right).
\]

The ring A' satisfies the hypotheses of the proposition, and its residue field is infinite. By part C) of the proof, we have \(\text{long}_{\text{A}'}(\text{B}'/\text{c}') \leqslant 2 \text{long}_{\text{A}'}(\text{B}'/\text{A}')\) , and the equality implies that A' is a Gorenstein ring; but this latter condition entails that A is a Gorenstein ring (§ 3, no. 8, cor. 1 of prop. 12).

\section{Local cohomology, Grothendieck duality}
